% !TeX program = xelatex
% -----------------------------------------------------------------------------
% This document is written to compile under EITHER XeLaTeX/LuaLaTeX OR
% pdfLaTeX: the iftex package below detects the engine actually running and
% loads the matching font/encoding setup automatically, so no manual
% compiler change is required.
%
% The line "% !TeX program = xelatex" above is a "magic comment" that
% Overleaf, TeXstudio, VS Code's LaTeX Workshop, etc. read to select the
% default compiler for this file. If your setup instead has an explicit
% compiler manually pinned to pdfLaTeX (in Overleaf: Menu -> Compiler), that
% manual setting can override this comment -- that is fine and expected: the
% detection below makes the document work correctly either way. XeLaTeX or
% LuaLaTeX is still the slightly preferred choice, since it renders
% Vietnamese diacritics (ễ, ả, ố, ...) through native Unicode font support
% rather than a legacy font-encoding workaround.
% -----------------------------------------------------------------------------
\documentclass[11pt,a4paper]{article}

% Page and typography settings from the Word template
\usepackage[
  top=1in,
  bottom=1in,
  left=1in,
  right=0.983in,
  headsep=0.25in,
  footskip=0.4in
]{geometry}

% Detect the compiler engine and load the matching Vietnamese-capable font
% setup: fontspec for XeLaTeX/LuaLaTeX, or fontenc+inputenc+tgtermes (with
% the vntex-provided T5 encoding) for pdfLaTeX.
\usepackage{iftex}
\ifPDFTeX
  \usepackage[T5,T1]{fontenc}
  \usepackage[utf8]{inputenc}
  \usepackage{tgtermes}
\else
  \usepackage{fontspec}
  % TeX Gyre Termes is a Times-like font with broad Unicode coverage,
  % including Vietnamese diacritics, and ships with every standard TeX
  % Live/Overleaf installation, so no extra font files need to be uploaded.
  \setmainfont{TeX Gyre Termes}
\fi
\usepackage{microtype}
\usepackage{setspace}

% Core packages
\usepackage{amsmath,amssymb}
\usepackage{graphicx}
\usepackage[table]{xcolor}
\usepackage{array,tabularx}
\usepackage{titlesec}
\usepackage{placeins}
\usepackage{hyperref}
\usepackage{tikz}
\usetikzlibrary{arrows.meta,positioning,shapes.geometric,fit,backgrounds}

% Template colours
\definecolor{headingred}{HTML}{C00000}
\definecolor{guidanceblue}{HTML}{0000FF}
\definecolor{tableheader}{HTML}{FFE8E1}

% General paragraph format: Calibri-like 11 pt, approximately Word's spacing
\setstretch{1.08}
\setlength{\parindent}{0pt}
\setlength{\parskip}{8pt}
\setlength{\emergencystretch}{2em}
\setcounter{secnumdepth}{0}
\setcounter{tocdepth}{2}
\hypersetup{hidelinks}

% Heading styles corresponding to Heading 1, 2, and 3 in the DOCX file
\titleformat{\section}
  {\color{headingred}\bfseries\fontsize{16}{19}\selectfont}
  {}{0pt}{}
\titleformat{\subsection}
  {\color{black}\bfseries\fontsize{13}{16}\selectfont}
  {}{0pt}{}
\titleformat{\subsubsection}
  {\color{black}\bfseries\fontsize{12}{14}\selectfont}
  {}{0pt}{}
\titlespacing*{\section}{0pt}{12pt}{2pt}
\titlespacing*{\subsection}{0pt}{8pt}{2pt}
\titlespacing*{\subsubsection}{0pt}{6pt}{2pt}

% Helpful template commands
\newcommand{\guidance}[1]{%
  {\color{guidanceblue}\itshape [#1]\par}%
}
\newcommand{\chapterstart}[1]{\clearpage\section{#1}}
\newcommand{\blankfield}{\rule{0pt}{2.7ex}\hspace{0.96\linewidth}}
% Under XeLaTeX/LuaLaTeX with a Unicode-capable font (set above), Vietnamese
% text needs no special encoding switch. Under pdfLaTeX, it must be typeset
% in the T5-encoded Vietnamese-capable font. Either way, the rest of the
% document just calls \vietnamesetext{...} and does not need to change.
\ifPDFTeX
  \newcommand{\vietnamesetext}[1]{{\fontencoding{T5}\fontfamily{qtm}\selectfont #1}}
\else
  \newcommand{\vietnamesetext}[1]{#1}
\fi

% Print generated lists without creating a second title
\makeatletter
\newcommand{\printtoc}{\@starttoc{toc}}
\newcommand{\printlistoftables}{\@starttoc{lot}}
\newcommand{\printlistoffigures}{\@starttoc{lof}}
\makeatother

% -----------------------------------------------------------------------------
% Project information: edit these values only
% -----------------------------------------------------------------------------
\newcommand{\ProjectName}{Intelligent Voice Customer Support Integrated with Retrieval-Augmented Generation}
\newcommand{\GroupName}{GSU26AI02}
\newcommand{\MemberOneName}{\vietnamesetext{Nguyễn Quốc Nhựt}}
\newcommand{\MemberTwoName}{\vietnamesetext{Nguyễn Võ Minh Đạt}}
\newcommand{\MemberThreeName}{\vietnamesetext{Hoàng Văn Lập}}
\newcommand{\MemberFourName}{\vietnamesetext{Nguyễn Quang Khải}}
\newcommand{\MemberOne}{\MemberOneName{} -- SE181753}
\newcommand{\MemberTwo}{\MemberTwoName{} -- QE180070}
\newcommand{\MemberThree}{\MemberThreeName{} -- SE180503}
\newcommand{\MemberFour}{\MemberFourName{} -- SE181512}
\newcommand{\SupervisorName}{\vietnamesetext{Nguyễn Quốc Trung}}
\newcommand{\ExternalSupervisorName}{\vietnamesetext{Nguyễn Xuân Huy}}
\newcommand{\CapstoneCode}{SU26AI47}
\newcommand{\SubmissionMonth}{[month]}
\newcommand{\SubmissionYear}{[year]}

\begin{document}

% -----------------------------------------------------------------------------
% Cover page
% -----------------------------------------------------------------------------
\hypersetup{pageanchor=false}
\begin{titlepage}
\thispagestyle{empty}

\begin{tabularx}{\textwidth}{@{}>{\centering\arraybackslash}p{0.38\textwidth}X@{}}
  \includegraphics[width=0.34\textwidth]{assets/fpt-education.png} &
  \centering\bfseries\fontsize{16}{19}\selectfont
  MINISTRY OF EDUCATION AND TRAINING
\end{tabularx}

\vspace{1.05cm}

\begin{center}
  {\bfseries\fontsize{36}{43}\selectfont FPT UNIVERSITY\par}

  \vspace{1.5cm}
  {\fontsize{28}{34}\selectfont Capstone Project Document\par}

  \vspace{0.55cm}
  {\fontsize{22}{27}\selectfont \ProjectName\par}

  \vspace{1.0cm}
  {\bfseries\fontsize{18}{22}\selectfont \GroupName\par}
\end{center}

\vspace{0.45cm}

\renewcommand{\arraystretch}{1.35}
\begin{tabularx}{0.94\textwidth}{@{}>{\raggedleft\arraybackslash\bfseries}p{0.31\textwidth}@{\hspace{1.2em}}X@{}}
  Group Members & \MemberOne \\
                & \MemberTwo \\
                & \MemberThree \\
                & \MemberFour \\
  Supervisor & \SupervisorName \\
  Ext Supervisor & \ExternalSupervisorName \\
  Capstone Project code & \CapstoneCode
\end{tabularx}

\vfill

\begin{center}
  {\fontsize{14}{17}\selectfont -- Ho Chi Minh City, August 2026 --}
\end{center}
\end{titlepage}
\hypersetup{pageanchor=true}
\pagenumbering{roman}

% -----------------------------------------------------------------------------
% Front matter
% -----------------------------------------------------------------------------
\clearpage
\section*{Table of Contents}
\printtoc

\clearpage
\section{Acknowledgement}
\guidance{Although not mandatory, a section thanking those who supported you
during the project, including advisors, lecturers, friends, and family, can be
included as an expression of gratitude.}

\section{Definition and Acronyms}
\guidance{A list of important terms and definitions used in the project. This
ensures that all readers have a common understanding of the specialized terms
used in the document.}

\renewcommand{\arraystretch}{1.35}
\begin{tabularx}{\textwidth}{|>{\centering\arraybackslash}p{0.15\textwidth}|X|}
\hline
\rowcolor{tableheader}
\textbf{Acronym} & \centering\arraybackslash\textbf{Definition} \\
\hline
AI  & Artificial Intelligence \\
\hline
DL  & Deep Learning \\
\hline
ML  & Machine Learning \\
\hline
PM  & Project Manager \\
\hline
PMP & Project Management Plan \\
\hline
WBS & Work Breakdown Structure \\
\hline
\ldots & \ldots \\
\hline
\blankfield & \\
\hline
\blankfield & \\
\hline
\end{tabularx}

\section{List of tables}
\guidance{A list of all the tables presented throughout the document.}
\printlistoftables

\section{List of figures}
\guidance{A list of all the figures presented throughout the document.}
\printlistoffigures

% -----------------------------------------------------------------------------
% I. Project Introduction
% -----------------------------------------------------------------------------
\clearpage
\pagenumbering{arabic}
\chapterstart{I. Project Introduction}

\subsection{1. Overview}

\subsubsection{1.1 Project Information}

\textit{\ProjectName} is a capstone project at FPT University with the project
code \CapstoneCode. The project is carried out by group \GroupName, which
consists of four members: \MemberOneName, \MemberTwoName, \MemberThreeName,
and \MemberFourName.

The project is developed under the academic supervision of \SupervisorName{}
and the external supervision of \ExternalSupervisorName. Its principal
stakeholders include the project team, the supervisors, FPT University,
organizations that provide customer-support services, customer-support agents,
and customers who interact with the proposed voice-based support system.

\subsubsection{1.2 Project Overview}

Our team develops an intelligent voice-based customer-support system for the construction domain. By combining spoken interaction with Retrieval-Augmented Generation (RAG), the system enables customers to ask questions naturally and receive relevant information obtained from an enterprise's internal knowledge sources. The generated response can be delivered as either text or synthesized speech, allowing users to select the most convenient form of interaction.

A key objective of the project is to improve speech recognition for Vietnamese
speakers from different regions. The proposed solution combines PhoWhisper with
RouteLoRA to process diverse regional accents more accurately. The recognized
request is then handled by an Agentic RAG pipeline, which identifies the user's
intent, selects appropriate information-retrieval actions, searches the
enterprise knowledge base, and generates a contextually grounded response.
This end-to-end approach aims to provide an accessible, accurate, and efficient
customer-support experience while reducing the workload associated with
repetitive inquiries.

\subsection{2. Project Background}

\textbf{Customer support in general and in the construction domain.}
Customer support includes the assistance an organization provides before,
during, and after a customer purchases or uses its products and services. In
many organizations, this work is centralized through call or contact centers,
which operate as socio-technical systems in which customer demand, employee
availability, service time, and information quality jointly affect the customer
experience \cite{gans2003}. Effective support is particularly important in the
construction domain because customers may need timely explanations of building
materials, product specifications, quotations, warranties, technical
procedures, project progress, and after-sales policies. Construction research
has identified communication, client orientation, responsiveness to complaints,
and service performance as important contributors to client satisfaction
\cite{ahmed1995}. Nevertheless, conversational-AI adoption in the architecture,
engineering, and construction (AEC) industry remains less mature than in many
other service sectors \cite{saka2023}.

\textbf{Traditional customer-support process.}
In a traditional process, a customer contacts an enterprise by telephone,
email, or an in-person channel. A human agent receives the request, determines
its category, searches manuals, frequently asked questions, product documents,
project records, or internal software, and may consult a technical department
before returning an answer. This workflow depends strongly on agent experience
and staffing capacity. When demand exceeds the number of available agents,
requests form queues and response times increase; differences in agent knowledge
can also produce inconsistent answers \cite{gans2003}. The problem is amplified
in construction because relevant information is often distributed across
technical documents and complex systems such as Building Information Modeling
(BIM). Nabavi et al. reported that manual BIM information retrieval can be
tedious and time-consuming for non-technical users and demonstrated that an
NLP-based inquiry platform could answer selected questions substantially faster
than manual expert search \cite{nabavi2023}.

\textbf{Need for digital transformation.}
Customers increasingly expect rapid, consistent, and continuously available
support across multiple communication channels. At the same time, enterprises
must control operating costs and allow human agents to focus on cases that
require judgment, negotiation, or empathy. Research on artificial intelligence
in service describes a progression from the automation of routine mechanical
tasks toward analytical and interactive service activities \cite{huang2018}.
AI-based chatbots have consequently become an important customer self-service
channel because they can provide immediate responses and reduce service time,
although their effectiveness still depends on dialogue quality, transparency,
and the ability to understand the user's actual request \cite{adam2021}. For the
construction sector, the limited maturity of conversational AI and the growing
volume of digital project and enterprise information create a clear need for a
more accessible natural-language interface \cite{saka2023}.

\textbf{Natural language processing as an enabling technology.}
Natural language processing (NLP) enables computers to analyze text or speech,
recognize user intent and important entities, retrieve relevant information,
and generate a response. A systematic review of NLP in customer service found
that chatbots and question-answering systems have been studied across multiple
service domains and support both textual and spoken interaction; it also
identified data volume, diversity, and quality as recurring limitations
\cite{mashaabi2022}. A broader review of service chatbots similarly shows the
transition from hand-crafted rules to machine-learning and deep-learning
approaches \cite{suhaili2021}. In construction, NLP has already enabled
non-expert users to query BIM information through ordinary language and voice
interfaces \cite{nabavi2023}. Retrieval-Augmented Generation (RAG) further
connects language generation to an external document collection, producing
answers that are more specific and factually grounded than those of a purely
parametric model \cite{lewis2020}. In customer-service question answering,
retrieval over structured organizational knowledge has also been shown to
improve answer quality and reduce issue-resolution time \cite{xu2024}.

\textbf{Previous conversational and construction-support systems.}
Early conversational programs such as ELIZA generated responses through
keyword-triggered decomposition and reassembly rules, creating the appearance
of dialogue without genuine domain understanding \cite{weizenbaum1966}. Later
customer-service chatbots introduced intent classification, learned language
representations, and more flexible dialogue mechanisms, but many systems still
remain restricted to predefined intents, narrow datasets, or text-only
interaction \cite{suhaili2021}. Within construction, BIMASR introduced a
voice-based framework for retrieving BIM information \cite{shin2021}; Lin
prototyped an NLU chatbot through which site managers could ask about building
component dimensions \cite{lin2023}; and Nabavi et al. developed an NLP-based
voice assistant for selected 4D BIM queries \cite{nabavi2023}. More recent work
has used prompt-based large language models to interpret natural-language BIM
queries and summarize retrieved results \cite{zheng2023}. These systems
demonstrate the feasibility of conversational access to construction data, but
they primarily target project practitioners and structured BIM use cases.
They do not fully address Vietnamese regional accents, general enterprise
documents, customer-facing conversations, and both text and voice output in a
single end-to-end system.

\textbf{Main motivation for the project.}
The project is motivated by the gap between the need for accessible,
enterprise-grounded construction support and the limitations of existing
systems. Vietnamese voice interaction introduces an additional difficulty
because a production system must recognize speakers with different regional
accents. PhoWhisper was fine-tuned on 844 hours of speech covering diverse
Vietnamese accents and achieved strong results on Vietnamese automatic speech
recognition benchmarks \cite{le2024}. Low-Rank Adaptation (LoRA) provides a
parameter-efficient method for specializing large pretrained models
\cite{hu2022}, while recent multi-accent ASR research shows that combining
accent-specific LoRA experts can improve recognition under both known- and
unknown-accent conditions \cite{bagat2025}. These findings motivate the
project's use of PhoWhisper with a RouteLoRA mechanism to select or combine
regional adaptation experts.

After speech recognition, the system must retrieve current information from
enterprise-controlled sources rather than rely only on knowledge stored in a
language model. RAG provides the grounding mechanism \cite{lewis2020}, while
the ReAct framework demonstrates how a language model can interleave reasoning
with actions that access external knowledge sources \cite{yao2023}. Combining
these ideas motivates an Agentic RAG pipeline that can analyze a customer's
request, choose appropriate retrieval actions, obtain construction-related
enterprise information, and generate a response as text or synthesized speech.
The overall motivation is therefore to deliver more accurate and inclusive
Vietnamese voice support, shorten response time, improve answer consistency,
and reduce the repetitive information-search workload of human agents.

\subsection{3. Project Objective}
\guidance{State the specific objectives of the project and what the project
hopes to achieve at the highest level. Objectives should be clear, measurable,
and achievable. They guide the development of the project and the evaluation
of its success.}

\subsection{4. Problem Statement}

Construction enterprises must respond to customer questions concerning
products, materials, technical specifications, quotations, warranties,
policies, and project-related procedures. The information required to answer
these questions is often distributed across manuals, internal documents,
project records, and specialized information systems. In the traditional
support process, human agents must interpret each request, search these sources,
and consult technical personnel when necessary. This process is difficult to
scale when the number of inquiries increases and can result in long waiting
times, repetitive workload, and inconsistent answers among agents
\cite{gans2003,nabavi2023}.

Existing conversational systems do not fully resolve this problem. Rule-based
and intent-based chatbots are generally limited to predefined scenarios, while
general-purpose language models may generate plausible responses that are not
supported by current enterprise information \cite{suhaili2021,lewis2020}. Prior
construction assistants have demonstrated natural-language or voice access to
BIM data, but they mainly support narrow queries for engineers, site managers,
or other project practitioners \cite{shin2021,lin2023,zheng2023}. Consequently,
there remains a gap for a customer-facing system that can search heterogeneous
construction-enterprise knowledge sources and provide grounded answers to a
broader range of customer questions.

Voice interaction introduces an additional challenge for Vietnamese users.
Differences among Northern, Central, and Southern accents, together with
speaking speed, environmental noise, and construction-specific terminology,
can cause automatic speech-recognition errors. An incorrectly transcribed
question may lead to incorrect intent interpretation, irrelevant retrieval,
and an unsuitable final response. Although PhoWhisper provides a strong
Vietnamese ASR foundation using speech data with diverse accents
\cite{le2024}, the system still requires an effective adaptation mechanism to
handle regional and domain variation consistently. Research on accent-specific
LoRA experts indicates that routing or combining lightweight adapters is a
promising direction for multi-accent speech recognition \cite{bagat2025}.

The central problem addressed by this project is therefore how to design an
end-to-end intelligent voice customer-support system for the construction
domain that can accurately recognize Vietnamese speech across regional
accents, retrieve reliable information from enterprise-controlled sources,
generate contextually grounded answers, and deliver those answers in either
text or synthesized speech. The system must integrate PhoWhisper and RouteLoRA
for accent-aware speech recognition with an Agentic RAG pipeline capable of
interpreting requests, selecting appropriate retrieval actions, and using
external knowledge during response generation \cite{lewis2020,yao2023}. In
addition to functional integration, the project must balance recognition
accuracy, answer relevance, response time, and usability so that the resulting
solution can provide practical support without removing the need for human
intervention in complex or sensitive cases.

\subsection{5. Significance of the Project}

The project has practical significance for both construction enterprises and
their customers. By providing a natural voice interface, the proposed system
can make enterprise information easier to access for users who may find typing
inconvenient or who need assistance while performing other activities. Faster
access to product specifications, policies, warranties, procedures, and other
construction-related information can reduce customer waiting time and improve
the consistency of the support experience. The option to receive the final
answer as either text or synthesized speech also allows the interaction channel
to be adapted to different user situations.

For enterprises, the system can automate a portion of repetitive information
requests and reduce the amount of time human agents spend searching across
multiple documents. This allows support personnel to focus on cases that
require technical expertise, negotiation, empathy, or direct coordination with
other departments. Grounding generated responses in enterprise-controlled
knowledge sources can also improve consistency and make it easier to update the
support service when products, policies, or internal documents change. These
benefits respond directly to the queueing, staffing, and information-access
challenges identified in traditional customer-support operations
\cite{gans2003,nabavi2023}.

The project is technically significant because it investigates two connected
problems that are often treated separately: Vietnamese multi-accent speech
recognition and enterprise-grounded conversational question answering.
PhoWhisper provides a Vietnamese ASR foundation trained on diverse accents
\cite{le2024}, while parameter-efficient LoRA experts provide a practical basis
for regional adaptation without maintaining a completely independent speech
model for every accent \cite{hu2022,bagat2025}. Integrating this accent-aware
speech pipeline with Agentic RAG extends the project beyond simple speech
transcription by enabling the system to reason about a request, retrieve
relevant evidence, and generate an answer grounded in external knowledge
\cite{lewis2020,yao2023}.

Within the construction industry, the project contributes to the ongoing
digital transformation from specialist-oriented information systems toward
more accessible conversational interfaces. Existing studies demonstrate the
potential of voice and natural-language access to BIM data, but customer-facing
applications and Vietnamese-language support remain limited
\cite{saka2023,shin2021,lin2023}. The proposed system therefore provides a
reference architecture and prototype for connecting construction-enterprise
knowledge with Vietnamese voice interaction. More broadly, the project may
support more inclusive access to information while preserving a clear path for
escalating complex, safety-critical, or sensitive requests to qualified human
personnel.

\subsection{6. Project Scope \& Limitations}

\textbf{Project scope.}
The project covers the design, implementation, and evaluation of a prototype
voice-based customer-support system for the construction domain. Its primary
users are Vietnamese-speaking customers seeking information from a selected
enterprise knowledge base. The supported content includes non-transactional
information such as product and material descriptions, technical
specifications, frequently asked questions, warranties, support policies, and
general procedures contained in documents supplied to the system.

The processing pipeline begins with a Vietnamese voice request. PhoWhisper is
used as the base automatic speech-recognition model, while RouteLoRA is used to
route or combine regional LoRA adapters intended to improve recognition across
Northern, Central, and Southern accents. The resulting transcript is passed to
an Agentic RAG component that interprets the request, selects an appropriate
retrieval action, searches the indexed enterprise documents, and provides
retrieved context to the response-generation model. The system then returns the
answer as text and may convert it into synthesized speech. The scope also
includes a basic interface through which users can submit requests and receive
responses.

Evaluation is limited to representative construction customer-support use
cases and the datasets and enterprise documents available during the project.
The speech component may be assessed using measures such as Word Error Rate or
Character Error Rate across regional-accent subsets. The retrieval and response
components may be evaluated through retrieval relevance, answer correctness,
faithfulness to the provided context, response latency, and selected usability
criteria. The project focuses on demonstrating the feasibility and integration
of the complete pipeline rather than replacing an enterprise's entire contact
center infrastructure.

\textbf{Project limitations.}
The speech-recognition performance depends on the quantity, quality, regional
coverage, and domain relevance of the available audio data. RouteLoRA can
improve adaptation, but it cannot guarantee accurate transcription for every
speaker, mixed accent, uncommon construction term, noisy environment, or
low-quality recording. The initial system is designed primarily for Vietnamese
and does not aim to provide equivalent performance for other languages.

The quality of Agentic RAG responses is constrained by the coverage, accuracy,
format, and freshness of the indexed enterprise documents. Information that is
missing, outdated, contradictory, or inaccessible cannot be reliably recovered.
Although retrieval grounding reduces unsupported generation, it does not
completely eliminate incorrect interpretation, irrelevant retrieval, or
hallucinated statements. The system must therefore provide uncertainty handling
and escalation to a human agent when sufficient evidence is unavailable or
when a request involves legal, financial, contractual, safety-critical, or
specialized engineering decisions.

Finally, the prototype is not intended to represent a fully production-ready
deployment. Large-scale concurrency, integration with every enterprise system,
continuous model monitoring, comprehensive cybersecurity certification, and
long-term operational maintenance are outside the immediate scope. Response
time and speech-output quality may also vary according to computing resources,
network conditions, model size, and third-party services. Consequently, results
obtained from the selected evaluation scenarios may not generalize directly to
all construction organizations, document collections, accents, or real-world
customer environments.

% -----------------------------------------------------------------------------
% II. Project Management Plan
% -----------------------------------------------------------------------------
\chapterstart{II. Project Management Plan}

The project is managed using a hybrid iterative approach. The major scope,
milestones, acceptance criteria, and academic deliverables are planned at the
beginning of the project, while technical development is divided into short
iterations. Each iteration produces a reviewable result such as a dataset
version, model baseline, retrieval pipeline, application feature, or evaluation
report. This approach allows the team to respond to experimental findings
without losing control of the overall schedule and project objectives.

\subsection{1. Team Work}

\subsubsection{1.1 Team Structure and Roles}

The team uses a functional structure in which each member owns a principal
technical area while remaining responsible for code review, integration,
testing, and documentation. The Project Manager coordinates dependencies and
serves as the primary contact for the supervisors. Technical decisions that
affect multiple components are discussed by the full team rather than being
made by a single component owner.

\begin{table}[htbp]
\centering
\caption{Team structure and primary responsibilities}
\label{tab:team-roles}
\small
\renewcommand{\arraystretch}{1.25}
\begin{tabularx}{\textwidth}{|p{0.22\textwidth}|p{0.22\textwidth}|X|}
\hline
\rowcolor{tableheader}
\textbf{Member} & \textbf{Primary role} & \textbf{Main responsibilities} \\
\hline
\MemberOneName & Project Manager and System Integration Lead &
Maintain the project plan and backlog; coordinate requirements, milestones,
meetings, and reports; define the overall architecture; manage cross-component
integration; monitor risks and changes; prepare demonstrations and submissions.
\\
\hline
\MemberTwoName & Speech AI and RouteLoRA Lead &
Prepare and analyze speech data; establish the PhoWhisper baseline; design and
train regional LoRA adapters; implement RouteLoRA; evaluate overall and
regional WER/CER; optimize inference and document ASR experiments. \\
\hline
\MemberThreeName & Knowledge Base and Agentic RAG Lead &
Prepare enterprise documents; design chunking, metadata, embedding, and
indexing processes; implement retrieval tools and agent workflows; evaluate
retrieval relevance, answer faithfulness, and source grounding. \\
\hline
\MemberFourName & Application, Quality Assurance, and Deployment Lead &
Develop the backend API and user interface; integrate text-to-speech and system
components; implement logging, error handling, automated tests, and deployment;
coordinate usability and end-to-end performance testing. \\
\hline
\SupervisorName & Academic Supervisor &
Provide academic and technical direction; review plans, methods, experiments,
reports, and milestone results; approve major scope or methodology changes. \\
\hline
\ExternalSupervisorName & External Supervisor &
Provide construction-domain and business guidance; review practical use cases,
enterprise information needs, and the applicability of the prototype. \\
\hline
\end{tabularx}
\end{table}
\FloatBarrier

Although the table assigns primary ownership, no component is developed in
isolation. Every merge into the main development branch requires review by at
least one other member. Each component owner must also provide setup notes,
interfaces, test cases, and progress information so that another member can
continue the work if necessary.

\subsubsection{1.2 Communication Plan}

The communication plan is designed to keep technical work synchronized and to
ensure that decisions can be traced later. Source code, issues, experiment
configurations, and change requests are maintained in a version-controlled
repository. Reports, meeting minutes, diagrams, datasets descriptions, and
evaluation results are stored in a shared document workspace. Short operational
updates are exchanged through a group messaging channel, while decisions that
affect scope, architecture, schedule, or evaluation must also be recorded in
the project documentation.

\begin{table}[htbp]
\centering
\caption{Project communication plan}
\label{tab:communication-plan}
\small
\renewcommand{\arraystretch}{1.25}
\begin{tabularx}{\textwidth}{|p{0.23\textwidth}|p{0.17\textwidth}|X|}
\hline
\rowcolor{tableheader}
\textbf{Communication activity} & \textbf{Frequency} &
\textbf{Purpose and expected record} \\
\hline
Asynchronous team update & Each working day &
Report completed work, the next task, and blockers. Significant blockers are
converted into tracked issues with an owner and expected resolution date. \\
\hline
Internal planning and review meeting & Weekly &
Review the backlog, demonstrations, experiment results, integration status,
risks, and work allocation. Meeting minutes record decisions and action items. \\
\hline
Component technical session & As required, normally weekly &
Resolve detailed ASR, RAG, application, data, or deployment issues. Design
decisions are recorded in technical notes or repository issues. \\
\hline
Academic supervisor meeting & Every two weeks and at milestones &
Present progress, receive methodological feedback, confirm academic quality,
and request approval for major changes. A concise progress report is prepared. \\
\hline
External supervisor review & At requirements, prototype, and final-review
milestones & Validate construction use cases, terminology, practical value,
and the suitability of enterprise-facing functions. \\
\hline
Urgent escalation & Within 24 hours of identification &
Notify the Project Manager immediately when a blocker threatens a milestone,
data access, security, or a core technical objective. \\
\hline
\end{tabularx}
\end{table}
\FloatBarrier

To avoid unresolved disagreements, the component owner first proposes a
technical recommendation supported by evidence. If consensus cannot be reached,
the Project Manager compares the alternatives against scope, quality, risk, and
schedule. Issues that materially affect project objectives or methodology are
escalated to the academic supervisor for a final decision.

\subsection{2. Project Management Approach}

Work is organized as a prioritized backlog and implemented in approximately
two-week iterations. Each task must have an owner, expected output, dependency,
acceptance condition, and target iteration. At the end of an iteration, the
team demonstrates completed work, reviews metrics and defects, and updates the
next iteration according to the remaining risks and milestone priorities.

\subsubsection{2.1 WBS}

The Work Breakdown Structure (WBS) divides the project into the following work
packages. The week ranges form an indicative fifteen-week plan and may be
adjusted to match the official capstone schedule without changing the required
deliverables.

\begin{enumerate}
  \item \textbf{Project initiation and requirements (Weeks 1--2).}
  Confirm stakeholders, customer-support use cases, functional and
  non-functional requirements, success criteria, constraints, communication
  procedures, and initial risks. The Project Manager leads this package with
  input from all members and both supervisors. Deliverables include the project
  charter, prioritized requirements, initial backlog, and management plan.

  \item \textbf{Literature review and solution architecture (Weeks 1--3).}
  Review Vietnamese ASR, parameter-efficient adaptation, multi-accent routing,
  RAG, agentic workflows, construction conversational systems, and evaluation
  methods. Define component boundaries, interfaces, data flow, technology
  choices, and experiment strategy. The Project Manager coordinates the work,
  while each technical lead reviews literature for the assigned component.

  \item \textbf{Data and knowledge preparation (Weeks 2--6).}
  Collect permitted speech resources and representative construction-domain
  enterprise documents; inspect licenses and sensitive content; clean,
  normalize, label, split, and version the data; define regional-accent subsets;
  and create the document ingestion, chunking, metadata, embedding, and indexing
  pipeline. The Speech AI Lead owns audio preparation, while the RAG Lead owns
  the enterprise knowledge base.

  \item \textbf{Speech-recognition development (Weeks 3--9).}
  Establish the PhoWhisper baseline, select LoRA target modules, train and
  compare regional adapters, implement the RouteLoRA mechanism, and evaluate
  WER/CER by region and test condition. Required resources include audio data,
  GPU compute, experiment tracking, model storage, and reproducible training
  configurations. The Speech AI Lead owns this package with peer review from
  the Project Manager.

  \item \textbf{Agentic RAG development (Weeks 4--10).}
  Implement document retrieval, metadata filtering, reranking where applicable,
  prompt templates, agent tools, source-grounded response generation, and
  insufficient-evidence handling. Construct test questions with reference
  documents and evaluate retrieval relevance, correctness, and faithfulness.
  The Knowledge Base and Agentic RAG Lead owns this package.

  \item \textbf{Application and voice-output development (Weeks 6--11).}
  Develop the user interface, backend services, audio upload or recording flow,
  session management, text response display, text-to-speech output, logging,
  and error messages. Define stable APIs connecting ASR, Agentic RAG, and TTS.
  The Application Lead owns this package with integration support from all
  component owners.

  \item \textbf{System integration and verification (Weeks 9--13).}
  Integrate all components into an end-to-end pipeline; implement unit,
  integration, regression, and failure-handling tests; measure latency and
  resource use; and correct interface or data-format defects. All members
  participate, coordinated by the Project Manager and Application Lead.

  \item \textbf{Evaluation and improvement (Weeks 11--14).}
  Execute controlled ASR, retrieval, generation, end-to-end, and usability
  evaluations. Compare results with defined baselines, analyze errors by
  accent and question category, prioritize improvements, and document remaining
  limitations. Component owners are responsible for their metrics, while the
  Project Manager consolidates the final evaluation.

  \item \textbf{Deployment, documentation, and closure (Weeks 13--15).}
  Prepare a reproducible prototype deployment, user and developer guides,
  architecture and API documentation, model and data cards, final report,
  demonstration, presentation, risk closure, lessons learned, and project
  archive. The Application Lead coordinates deployment and the Project Manager
  coordinates the remaining closure artifacts.
\end{enumerate}

\subsubsection{2.2 Risk Management}

Risks are reviewed during the weekly team meeting and before every milestone.
Each risk has an owner, likelihood, impact, mitigation action, and trigger for
contingency action. Likelihood and impact are classified as Low (L), Medium
(M), or High (H). A high-impact risk is escalated immediately when its trigger
is observed.

\begin{table}[htbp]
\centering
\caption{Major project risks and response strategies}
\label{tab:risk-register}
\footnotesize
\renewcommand{\arraystretch}{1.2}
\begin{tabularx}{\textwidth}{|p{0.22\textwidth}|>{\centering\arraybackslash}p{0.07\textwidth}|>{\centering\arraybackslash}p{0.07\textwidth}|X|}
\hline
\rowcolor{tableheader}
\textbf{Risk} & \textbf{L} & \textbf{I} & \textbf{Mitigation and contingency} \\
\hline
Insufficient or imbalanced regional speech data & H & H &
Audit accent coverage early; use stratified splits and augmentation; prioritize
the most underrepresented region; report per-region limitations rather than
claiming unsupported generalization. \\
\hline
ASR accuracy remains below the baseline or target & M & H &
Maintain a reproducible PhoWhisper baseline; run controlled adapter experiments;
analyze errors by accent, noise, and terminology; simplify routing or fall back
to the best validated adapter if RouteLoRA is unstable. \\
\hline
Enterprise documents are unavailable, outdated, or inconsistent & M & H &
Define the minimum document set and ownership early; record provenance and
version; create a representative approved test corpus; return an
insufficient-evidence response for unsupported questions. \\
\hline
RAG retrieves irrelevant evidence or generates unsupported answers & M & H &
Use metadata, retrieval evaluation, reranking where justified, source display,
faithfulness tests, confidence rules, and human escalation for sensitive cases. \\
\hline
End-to-end latency or compute cost is excessive & M & M &
Profile each component; cache reusable data; use smaller or quantized models
when quality remains acceptable; apply asynchronous processing; define a
degraded text-only mode if speech generation becomes a bottleneck. \\
\hline
Component integration fails near a milestone & M & H &
Define APIs and schemas early; provide mocks; integrate incrementally in every
iteration; maintain automated smoke tests and a known working release. \\
\hline
Schedule delay or temporary member unavailability & M & M &
Keep tasks small and documented; require cross-review; maintain schedule buffer;
reassign critical tasks and reduce low-priority features before reducing core
quality or evaluation activities. \\
\hline
Privacy, security, or licensing issue & M & H &
Use only authorized data; remove personal or confidential information; control
access and secrets; document licenses; avoid logging raw sensitive audio; stop
the affected experiment until approval is obtained. \\
\hline
Third-party model, API, or service becomes unavailable & L & M &
Abstract external services behind interfaces; record versions and
configurations; retain a local or alternative provider for core demonstrations
where feasible. \\
\hline
\end{tabularx}
\end{table}
\FloatBarrier

\subsubsection{2.3 Quality Management}

Quality management applies to requirements, data, models, software, evaluation,
and documentation. A requirement is considered complete only when its
acceptance condition is demonstrated and its evidence is stored. The team uses
version control, task tracking, peer review, reproducible configurations, and
release tags to prevent undocumented changes.

\begin{itemize}
  \item \textbf{Requirements quality:} maintain traceability from each
  high-priority requirement to its design component, implementation task, test,
  and final result. Ambiguous requirements are reviewed with the relevant
  supervisor before implementation.

  \item \textbf{Data quality:} document sources, permissions, preprocessing,
  labels, regional distribution, train-validation-test separation, known bias,
  and version. Duplicate or leaked evaluation samples must be removed.

  \item \textbf{ASR quality:} compare PhoWhisper, individual LoRA adapters, and
  RouteLoRA under the same data splits. Report WER/CER overall and separately
  for Northern, Central, and Southern speech, together with error analysis for
  noise and domain terminology.

  \item \textbf{RAG quality:} maintain a test set of questions, relevant source
  passages, and expected answer facts. Measure retrieval relevance and evaluate
  correctness, completeness, faithfulness, citation or source support, and the
  ability to refuse unsupported questions.

  \item \textbf{Software quality:} use peer-reviewed changes, automated unit
  tests for core functions, integration tests for component interfaces,
  end-to-end regression scenarios, structured logging, and clear error handling.
  Critical defects block a release until they are corrected or an approved
  workaround is documented.

  \item \textbf{Performance and usability quality:} report end-to-end latency
  and component-level timing using consistent test conditions. Conduct
  representative task-based user evaluation for clarity, ease of use, perceived
  response quality, and the text/speech interaction flow.

  \item \textbf{Documentation quality:} keep installation, configuration,
  architecture, API, experiment, evaluation, security, and user documentation
  synchronized with each milestone release. Report limitations and negative
  results rather than omitting them.
\end{itemize}

Before a milestone is accepted, the responsible owner performs self-testing,
another member performs a review, and the integrated result is demonstrated to
the team. Major academic deliverables are additionally reviewed by the
supervisor. This layered review process separates component completion from
actual project acceptance.

\subsubsection{2.4 Budget and Funding (Optional)}

The project is designed to minimize direct expenditure by prioritizing
open-source models, existing personal or university development equipment, and
free academic service tiers where permitted. However, a budget register is
maintained because model training, hosted inference, storage, data preparation,
and deployment may create variable costs.

\begin{table}[htbp]
\centering
\caption{Budget categories and financial controls}
\label{tab:budget-plan}
\small
\renewcommand{\arraystretch}{1.25}
\begin{tabularx}{\textwidth}{|p{0.24\textwidth}|p{0.27\textwidth}|X|}
\hline
\rowcolor{tableheader}
\textbf{Cost category} & \textbf{Planned approach} & \textbf{Control measure} \\
\hline
GPU training and inference & Use available local or academic compute first;
rent cloud GPU time only for experiments that cannot be completed locally. &
Set time and spending limits; terminate idle resources; record the cost of each
experiment configuration. \\
\hline
LLM, embedding, or TTS services & Prefer local or open-source alternatives;
use paid APIs only when they provide a necessary capability or comparison. &
Centralize credentials, set usage quotas, monitor requests, and require team
approval before increasing a paid limit. \\
\hline
Storage and deployment & Store only required datasets and model artifacts;
use a small prototype hosting environment. &
Apply retention rules, compress or archive inactive artifacts, and separate
development from demonstration resources. \\
\hline
Data labeling and domain validation & Distribute internal labeling tasks and
request expert review only for high-value domain cases. &
Use labeling guidelines, sampled quality checks, and prior approval for any
external paid work. \\
\hline
Contingency & Reserve capacity for unexpected compute, storage, or service
needs near evaluation and demonstration milestones. &
Use contingency only after documenting the reason, expected benefit, and impact
of not spending it. \\
\hline
\end{tabularx}
\end{table}
\FloatBarrier

The Project Manager records estimated and actual costs, the responsible member,
and the associated work package. Any expense beyond the agreed project limit
must be reviewed by the team and, where required, approved by the supervisor.
Financial decisions must not compromise data privacy, licensing, or evaluation
integrity.

\subsubsection{2.5 Change Management Process}

Changes may originate from supervisor feedback, new user requirements,
experimental results, unavailable data, technical constraints, risks, or defect
reports. All non-trivial changes are handled through the following controlled
process:

\begin{enumerate}
  \item \textbf{Submit and record the request.} The requester describes the
  proposed change, rationale, affected requirement or component, urgency, and
  expected benefit in the change log.

  \item \textbf{Perform impact analysis.} The relevant component owners estimate
  effects on scope, architecture, data, security, quality, schedule, budget,
  documentation, and other components. Alternatives, including rejecting or
  postponing the change, are documented.

  \item \textbf{Classify and decide.} A minor change does not alter a core
  objective, milestone, public interface, or evaluation method and may be
  approved by the Project Manager with team agreement. A major change affects
  one of these areas and requires supervisor approval.

  \item \textbf{Plan and implement.} After approval, the backlog, owner,
  priority, acceptance criteria, schedule, risks, and documentation are updated.
  Implementation follows the normal review and testing process.

  \item \textbf{Verify and close.} The requester or designated reviewer confirms
  that the acceptance criteria are satisfied and that related documents and
  tests have been updated. The final decision and release containing the change
  are recorded.
\end{enumerate}

Emergency fixes for security or complete service failure may be implemented
immediately, but they must be documented, reviewed, and tested retrospectively
before the next release. Unapproved scope expansion is not included merely
because implementation has already started.

\subsubsection{2.6 Closure and Evaluation}

Project closure begins with a feature freeze before the final evaluation. Only
defect corrections, documentation updates, and explicitly approved critical
changes are permitted after that point. The team then executes the final test
plan on a tagged release and confirms that high-priority requirements have
evidence of completion or are clearly recorded as limitations.

The final evaluation consolidates ASR results by regional accent, retrieval and
answer-quality results, end-to-end latency, resource use, failure cases, and
usability observations. Results are compared with the defined baselines and
acceptance criteria. The team also reviews whether the project remained within
its approved scope, schedule, cost controls, quality process, and ethical and
security constraints.

Closure deliverables include the final source code, dependency and deployment
instructions, model configurations, permitted data descriptions, knowledge-base
ingestion procedure, architecture and API documentation, test cases, evaluation
report, user guide, final report, presentation, and demonstration materials.
Secrets, temporary data, obsolete artifacts, and unauthorized copies are
removed before the project archive is created.

Finally, the team conducts a retrospective meeting with the following agenda:
what worked well, what did not work, which assumptions were incorrect, which
risks occurred, how communication and estimation could improve, and which
technical directions should be continued as future work. Lessons learned and
remaining risks are documented, unresolved issues are assigned a disposition,
and the prototype and documentation are handed over according to the
supervisors' requirements. The project is considered closed only after final
acceptance, archive verification, and completion of all required academic
submissions.

% -----------------------------------------------------------------------------
% III. Existing Systems / State of the Art
% -----------------------------------------------------------------------------
\chapterstart{III. Existing Systems/State of the Art}
\guidance{Provide the final existing-systems information following the
template as part II in Report 3.}
\guidance{This section is usually titled ``Literature Review'' or ``Related
Work.'' For projects related to specific technologies or systems, it may be
titled ``Existing Systems'' or ``State of the Art.'' Choose a title that fits
the content and context of the project.}

\subsection{1. Overview of the Field}

\subsection{2. Historical Context}
\guidance{Provide historical background or developments in the field,
especially those directly related to the project.}

\subsection{3. Key Studies and Theories}
\guidance{Describe and assess key studies, theories, and published works
related to the topic. Include both fundamental and applied research.}

\subsection{4. Technological Advancements}
\guidance{Discuss recent or significant technological advancements in the
field and how they affect the project.}

\subsection{5. Comparison of Existing Systems}
\guidance{If the project involves systems or technology, compare existing
systems or technologies, including their advantages and disadvantages in
relation to the proposed project.}

\subsection{6. Gaps in the Literature/Technology}
\guidance{Identify gaps or limitations in current research or technological
systems and explain how the project intends to address or improve upon them.}

\subsection{7. Justification for the Project}
\guidance{Based on the information above, highlight the necessity and
uniqueness of the project. Explain why the approach to the problem is important
and how it differs from or improves upon existing research or systems.}
\guidance{This section is not just a list of studies; it is an analysis and
evaluation that connects existing systems, research, and related work to the
project and shows their impact on it.}

% -----------------------------------------------------------------------------
% IV. Methodology
% -----------------------------------------------------------------------------
\chapterstart{IV. Methodology}

The project follows two coordinated methodological tracks that converge into a
single voice-based customer-support pipeline. The first track develops and
evaluates Vietnamese multi-accent speech recognition using PhoWhisper and
RouteLoRA. The second track develops and evaluates an Agentic Retrieval-
Augmented Generation (RAG) system that answers customer questions grounded in
enterprise construction-material documents, with a particular focus on
long, multi-page pricing tables. This chapter documents the research
questions, data, compared configurations, metrics, and quality-control
procedures for both tracks. The Agentic RAG methodology below reflects
completed experimental work; the ASR methodology is presented as a
structured outline to be completed by the Speech AI Lead.

\subsection{1. Research Questions and Objectives}

The project's overall objective is to determine, for each pipeline stage,
which design choices measurably improve the accuracy, reliability, and
auditability of a Vietnamese voice-and-text customer-support system, and at
what cost. This objective is pursued through two families of research
questions.

\textbf{ASR track.} \guidance{State the specific research questions for the
PhoWhisper/RouteLoRA multi-accent recognition study, e.g., how does regional
accent affect baseline WER/CER, and does RouteLoRA improve recognition over a
single fine-tuned model or single-accent adapters, under known- and
unknown-accent conditions.}

\textbf{Agentic RAG track.} Because a large share of the enterprise knowledge
base consists of long tables (construction-material price bulletins in which
a single table can span several hundred pages without repeating its header
row), the central research question is:

\begin{quote}
\itshape How does processing and splitting documents according to table
structure, as opposed to splitting flattened plain text by length, affect
structure preservation, embedding quality, evidence retrieval, and answer
accuracy; for which question types does the benefit appear, and at what cost?
\end{quote}

To avoid attributing every difference to a single end-metric, the system is
evaluated as four dependent layers: structure preservation
$\rightarrow$ discriminability in embedding space $\rightarrow$ evidence
retrieval $\rightarrow$ the ability to use retrieved context to generate an
answer. This layered design makes it possible to determine whether a
table-aware pipeline helps primarily at the retrieval step or primarily in
how the generation model interprets evidence once it is already in context.
The track answers six research questions:

\begin{itemize}
  \item \textbf{RQ1 -- Structure preservation.} Does a table-aware pipeline
  preserve column labels and row--column relationships better than recursive
  chunking?
  \item \textbf{RQ2 -- Representation cost.} Does repeating the header row
  make chunks more similar to one another in embedding space, reducing a
  dense retriever's ability to discriminate between them?
  \item \textbf{RQ3 -- Retrieval.} How do table-aware and recursive chunking
  differ in Recall@1/3/5/10, MRR, and the top-5 operating point; and, holding
  chunk boundaries fixed, how does the string representation used for
  embedding (HTML, key--value, or verbalized prose) affect dense retrieval?
  \item \textbf{RQ4 -- Answer generation.} Does the table-aware pipeline
  improve Exact Match, particularly for questions that require mapping a
  value to the correct column?
  \item \textbf{RQ5 -- Context-use mechanism.} When retrieval is not better
  but generation is, does the evidence support the hypothesis that the
  table-aware pipeline converts retrieved evidence into a correct answer more
  effectively?
  \item \textbf{RQ6 -- Robustness and reliability.} Do the conclusions hold
  when lexical (BM25) retrieval is added, and when the system is evaluated on
  questions that have no answer in the corpus?
\end{itemize}

\subsection{2. ASR Methodology: PhoWhisper and RouteLoRA}
\guidance{To be completed by the Speech AI Lead, following the same
evidence-based style as Section 3 below. Suggested structure:}

\subsubsection{2.1 Speech Data Collection and Preprocessing}
\guidance{Describe sources of Vietnamese speech data, regional-accent
coverage (Northern/Central/Southern), licensing, cleaning, normalization, and
train/validation/test data splits.}

\subsubsection{2.2 PhoWhisper Baseline}
\guidance{Describe the base PhoWhisper checkpoint used, fine-tuning or
zero-shot configuration, and baseline WER/CER overall and by accent.}

\subsubsection{2.3 Regional LoRA Adapter Design}
\guidance{Describe LoRA target modules, rank, and training configuration for
each regional adapter, and the data used to train each one.}

\subsubsection{2.4 RouteLoRA Mechanism}
\guidance{Describe how the routing or combination mechanism selects or mixes
regional adapters at inference time, and any routing-confidence or
fallback behaviour.}

\subsubsection{2.5 ASR Evaluation Metrics and Statistical Procedure}
\guidance{Describe WER/CER computation, per-accent breakdown, significance
testing between configurations, and reproducibility controls (fixed test
sets, fixed decoding parameters, logged model/checkpoint versions).}

\subsection{3. Agentic RAG Methodology}

\subsubsection{3.1 Evaluation Framework and Scope}

The unit of evaluation is a complete, table-aware, end-to-end pipeline:
PDF $\rightarrow$ table-region detection and extraction $\rightarrow$
cell-grid reconstruction $\rightarrow$ HTML table representation
$\rightarrow$ header identification or recovery $\rightarrow$ splitting at
row boundaries $\rightarrow$ header repetition per chunk $\rightarrow$
vector encoding $\rightarrow$ retrieval $\rightarrow$ answer generation. The
recursive baseline is likewise a complete pipeline, starting from flattened
PDF text split by length. Reported results therefore reflect the effect of
the entire structure-aware pipeline, not only the chunk-boundary rule in
isolation.

The chapter's principal comparison is denoted \textbf{T1500--R1500}: both
branches use the same nominal chunk-size limit of 1,500 tokens, the same
question set, the same embedding model, the same top-$k$, and the same
generation model. Holding these factors fixed isolates the effect of table
representation and chunk boundaries from other confounds. A secondary
configuration, \textbf{T3000} (table-aware, 3,000-token limit), is retained
only as a reference point for the effect of chunk size on chunk count and
embedding similarity, and is not substituted for the T1500--R1500 comparison.

\subsubsection{3.2 Document Corpus and Question Sets}

The corpus consists of 11 PDF documents -- official price bulletins and
their annexes for construction materials -- from which a structured
extractor produced 11,525 price rows used to build questions and analyze
structure. The largest document contains a single table that runs
continuously for approximately 699 pages with a header that appears only at
the start, making it the case where a header-repetition strategy has the
strongest basis for making a measurable difference. Five Markdown files that
had previously been part of the data pool were excluded from the benchmark
because they pass through a different pipeline and do not produce HTML table
chunks; keeping them would dilute measurement of the components under
evaluation.

\begin{table}[htbp]
\centering
\caption{Evaluation question sets}
\label{tab:rag-question-sets}
\small
\renewcommand{\arraystretch}{1.25}
\begin{tabularx}{\textwidth}{|X|>{\centering\arraybackslash}p{0.12\textwidth}|X|}
\hline
\rowcolor{tableheader}
\textbf{Set} & \textbf{Count} & \textbf{What it measures} \\
\hline
Price lookup (G1--G4): full name, name without diacritics, character-split
product code, shortened name & 252 &
Finding the single correct numeric value attached to a material name or code \\
\hline
Table attribute (S1--S4): unit, manufacturer, technical standard, price basis
& 248 &
Identifying the correct value in the correct column of the correct row \\
\hline
\textbf{Total answerable} & \textbf{500} & \\
\hline
Unanswerable (N1: existing row, requested cell empty -- 30; N3: near-miss
name/specification that does not exist -- 20; N4: requested publication
period absent from the corpus -- 20) & 70 &
Whether the system correctly refuses instead of fabricating a value \\
\hline
\end{tabularx}
\end{table}
\FloatBarrier

Splitting the 500-question benchmark into price-lookup and attribute groups
is a deliberate design choice. A price-lookup question can be answered
correctly whenever a product name and a plausible number appear near each
other, even without full structural interpretation, whereas an attribute
question forces the system to identify a string belonging to the correct
column of the correct row. Every result table therefore reports both the
overall figure and the two group-level figures. On the 70-question
unanswerable set, a response is scored correct only when the system declines
to give a specific value and states that the information could not be found;
this set is not interpreted independently of accuracy on the 500 answerable
questions, since a system that refuses every query could show a low
fabrication rate while being practically useless.

\subsubsection{3.3 Compared Configurations}

\textbf{Table-aware pipeline.} Table geometry is reconstructed with
\texttt{pdfplumber}, tables are represented as HTML, headers are recovered
for continuation pages, and tables are split at row boundaries; whenever a
table produces multiple chunks, the header row is repeated in each chunk and
counted against the token budget. Two chunk-size configurations are used:
T1500 (1,500-token limit) for the main generation comparison, and T3000
(3,000-token limit) as a reference point.

\textbf{Recursive baseline.} A \texttt{RecursiveCharacterTextSplitter} is
applied to flattened PDF text with no cell markup retained. The main
comparison configuration, R1500, is subject to the same nominal 1,500-token
limit as T1500 but produces fewer chunks because it neither repeats headers
nor preserves HTML structure.

\textbf{Retrieval configurations.} Dense retrieval uses
\texttt{text-embedding-3-small}. Hybrid retrieval fuses dense retrieval with
BM25~Okapi using Reciprocal Rank Fusion (RRF, $k=60$); each branch prefetches
50 candidates before fusion, and the top~5 fused results form the context for
generation. Applying BM25 to both T1500 and R1500 produces a $2\times2$
design (chunking $\times$ retriever) that separates two questions: whether
table-aware chunking outperforms recursive chunking under the same
retriever, and whether BM25 helps table-aware chunking more than it helps
recursive chunking or is simply a general retrieval improvement.

A separate \textbf{representation experiment} isolates the effect of string
format from the effect of chunk boundaries: three indices --
\texttt{T1500\_HTML}, \texttt{T1500\_KV}, and \texttt{T1500\_VERB} -- share
the same 2,252 logical chunks, rows, columns, headers, and cell values;
only the string passed to the embedding model changes. Key--value and
verbalized strings are produced by deterministic rules, not by a language
model. All three use the same embedding model, cosine similarity, and dense
retrieval, and this experiment does not call a generation model.

\textbf{Generation setup.} The primary generation model is
\texttt{openai/gpt-oss-20b}, run at \texttt{temperature=0},
\texttt{max\_tokens=2000}, top-$k=5$ context chunks, with tool calling
disabled, using the same prompt and answer-normalization function across all
four main configurations. All benchmark generation is executed directly
through OpenRouter with the same model name and parameters; results from any
other inference endpoint are never substituted in.

\subsubsection{3.4 Evaluation Metrics and Statistical Procedure}

\textbf{Structure layer.} Computed over all 11,525 rows without invoking a
language model: row integrity (material name and price remain in the same
chunk), rows with column labels (the chunk carries the header needed to
interpret the row), column-mapping capability, and rows excluded for
exceeding the chunk limit.

\textbf{Embedding layer.} Cosine similarity is computed between chunk pairs
within the corpus, reporting the median, the 90th percentile, and the share
of pairs with cosine similarity at or above 0.90 and 0.95. A high share of
near-duplicate pairs indicates an index with many vectors that are difficult
to discriminate, which can crowd several top-$k$ positions with near-duplicate
chunks.

\textbf{Retrieval layer.} Retrieval is evaluated at several cutoffs rather
than a single $k$: Recall@1/3/5/10 (the share of questions for which the
top-$k$ contains at least one chunk carrying the target value from
\texttt{expect\_values} or \texttt{expect\_text}), Mean Reciprocal Rank
(MRR), median first-hit rank, and No-hit@10. Recall@5 remains the
\textbf{primary operating-point metric} because the system's generation
step uses the top-5 context; Recall@1/3/10 and MRR are reported to describe
the shape of the ranking and avoid conclusions drawn from a single cutoff.
Because the benchmark does not yet carry an independent gold row or cell ID
for all 500 questions, these recall figures are \emph{proxy evidence hits}
(a query counts as a hit when the top-$k$ contains the expected value) --
useful for comparing configurations on the same benchmark, but not strictly
equivalent to evidence recall based on the correct row or cell; this
limitation is kept explicit throughout the interpretation of results.

\textbf{Generation layer.} Answers are scored by Exact Match (EM) after
normalization, with a denominator of 500 for every configuration. A
supplementary metric, \textbf{retrieval-to-answer conversion}, is defined
per \texttt{query\_id} as

\[
\text{Conversion}=
\frac{\text{questions with a retrieval hit \emph{and} a correct answer}}
{\text{questions with a retrieval hit}},
\]

which measures the ability to turn a retrieval hit into a correct answer,
while remaining descriptive because the underlying retrieval hit is itself
an answer-containing proxy.

\textbf{Rejection ability.} On the 70 unanswerable questions, the primary
metric is the correct-rejection rate; a response is scored incorrect when
the model states a specific value that is not supported by the documents.

\textbf{Statistical testing.} Because every configuration answers the same
question set, correct/incorrect comparisons use the paired McNemar test,
reporting both discordant counts (T correct/R incorrect and T
incorrect/R correct). The primary confirmatory family comprises six
T1500--R1500 comparisons (dense and hybrid, each on the full set, the
price-lookup subset, and the attribute subset); the six resulting p-values
are adjusted with the Holm procedure, which is more conservative than
correcting each subgroup of four separately and avoids selecting a testing
family after observing the results. BM25 interaction is measured as a
difference-of-differences,

\[
\Delta_{\text{interaction}}=
(T_{\text{hybrid}}-T_{\text{dense}})-(R_{\text{hybrid}}-R_{\text{dense}}),
\]

with uncertainty estimated by paired bootstrap over \texttt{query\_id}
(100,000 resamples, seed \texttt{20260807}, 95\% percentile confidence
interval for the generation interaction).

\subsubsection{3.5 Quality Control and Reproducibility}

\textbf{Index identity verification.} An earlier configuration alias for
``T'' silently defaulted to T3000, creating a risk of misattributing T3000
results to T1500. Evaluation scripts subsequently require explicit,
unambiguous configuration names (\texttt{T1500}, \texttt{T3000},
\texttt{R1500}); ambiguous aliases are rejected. Every vector cache carries a
manifest recording the configuration name, chunk limit, chunk count, a
content digest, the embedding model, and a cache version, and a cache is
reused only when the chunk count and digest match the current data.

\textbf{Raw-cache and aggregate separation.} Raw answer caches and
aggregated result files are written to separate paths. Each raw record
stores, at minimum, the \texttt{query\_id}, configuration, retrieved chunk
list, raw and normalized answers, correctness score, finish reason, an
empty-output flag, token usage, and the prompt hash and model, so that
aggregate files can never overwrite per-question data and McNemar,
conversion, and bootstrap statistics can be recomputed directly from raw
results.

\textbf{Endpoint consistency.} All official \texttt{openai/gpt-oss-20b}
results are generated directly through OpenRouter; the four main
configurations share the same prompt, normalization function, generation
parameters, and endpoint, and results from a second endpoint are never
substituted in.

\textbf{Label--parser dependency.} Expected answers are derived from
structured data produced by the same table-extraction pipeline used for the
table-aware branch, which could bias results in its favor. For the
row-integrity metric, only names that appear in every branch were retained
for comparison, after which T and R1500 were nearly identical. A random
sample of 25 attribute questions was cross-checked manually against the
source PDFs; 24 matched (96\% in this small sample), a result that reduces
but does not eliminate concern about systematic labeling error.

\textbf{Retrieval-label limitations.} The retrieval benchmark currently uses
\texttt{expect\_values} and \texttt{expect\_text} rather than an independent
gold row/cell ID. Because the same unit, manufacturer, or number can occur in
multiple rows, Recall@5 may contain false hits, and conversely a correct
generation can occur even when the retrieval proxy does not register a hit.
Conclusions about mechanism drawn from conversion are therefore presented as
consistent evidence rather than a definitive causal decomposition.

\subsubsection{3.6 Ethical and Data Considerations}

All source documents are publicly issued government or supplier price
bulletins for construction materials and contain no personal data. Vector
caches, raw evaluation logs, and configuration manifests are stored without
access credentials or other sensitive content. Because the system is
designed to surface a specific price to an end customer, the evaluation
protocol treats an incorrect but confident numeric answer as a materially
more harmful failure mode than a correct refusal, which motivates the
explicit measurement of rejection ability (Section 3.4) as a first-class
metric rather than a secondary concern.

% -----------------------------------------------------------------------------
% V. System Design and Implementation
% -----------------------------------------------------------------------------
\chapterstart{V. System Design and Implementation}

This chapter describes the design of the voice-based customer-support system,
cross-referenced against the running implementation rather than a purely
theoretical design. Design decisions that are backed by the benchmark
evidence in Chapter VI are identified as such; decisions that follow from a
design principle without a dedicated quantitative benchmark, and constraints
imposed by the deployment environment, are also identified explicitly so
that the two are not confused.

\subsection{1. Overall System Architecture}

\subsubsection{1.1 Business Context and Request Types}

Construction-material enterprises must draw information from heterogeneous
sources: government price bulletins, multi-page price annexes, supplier
quotations, and technical-standard documents. Much of the important data is
not continuous prose but resides in tables with the following
characteristics: a single table may span from tens to hundreds of pages;
header rows may appear only on the first page or may not repeat reliably;
many columns use vertically merged cells that can span an entire product
family; product name, specification, manufacturer, unit, and price sit in
different columns; some PDFs carry a text layer while others are scanned
images; the same product can carry several prices depending on period,
region, manufacturer, or price basis; and customer questions are typically
short, non-standardized, and reference near-identical product names.

The system must serve at least two categories of request that cannot be
served by the same mechanism, summarized in Table~\ref{tab:request-types}.

\begin{table}[htbp]
\centering
\caption{Request categories and the mechanism each one requires}
\label{tab:request-types}
\small
\renewcommand{\arraystretch}{1.3}
\begin{tabularx}{\textwidth}{|p{0.22\textwidth}|X|p{0.22\textwidth}|}
\hline
\rowcolor{tableheader}
\textbf{Category} & \textbf{Example} & \textbf{Appropriate mechanism} \\
\hline
Exact numeric lookup &
``What is the price of PCB40 cement in Hanoi?'' -- must match the exact
product, unit, period, and price; no ``approximate'' value is acceptable &
Structured-data query through a typed tool and SQL \\
\hline
Structural understanding &
``Who manufactures this product?'', ``Which standard applies to this
line?'' -- must correctly attribute a value to its column and row &
Table-aware RAG \\
\hline
Mixed request &
``What is the price of this product, and what is the applicable technical
standard?'' -- needs both an exact number and a document-grounded
explanation &
Tool + table-aware RAG combined \\
\hline
\end{tabularx}
\end{table}
\FloatBarrier

An architecture built on RAG alone cannot guarantee that a numeric answer is
attached to the correct row (Section 3.2 gives quantitative evidence for
this); an architecture built on structured SQL alone cannot handle ambiguous
phrasing, free-text conditions, notes, standards, and semantic relationships
within a table. The system is therefore designed as \textbf{two
complementary data paths} that share a single table-extraction layer at
ingestion time.

\subsubsection{1.2 Design Goals and Guiding Principles}

The architecture targets five goals: (1)~\textbf{preserve long-table
structure} -- retain headers, merged cells, and row--column relationships
across pages; (2)~\textbf{auditable numeric lookup} -- every price must trace
from its source cell into the database and be queried deterministically;
(3)~\textbf{support natural-language questions} -- users should not need to
know column names, schema, or SQL syntax; (4)~\textbf{fit common enterprise
infrastructure} -- prioritize CPU-only processing and avoid a hard
dependency on GPU or on running every PDF through visual OCR; and
(5)~\textbf{explainability and reproducibility} -- every answer must trace
back to its source document, chunk, or price row, and every design decision
should distinguish clearly between what was measured, what follows from a
stated principle, and what is an environment-imposed constraint.

Two principles run through the entire design:

\begin{quote}
\textbf{Principle 1 -- Numbers travel a deterministic path; text travels an
approximate-retrieval path.} Unit prices, quantities, and arithmetic are
always taken from structured data or a code-level calculation, never from a
language model. Conditions of applicability, standards, specifications, and
explanations are retrieved through RAG. This is not merely a design
preference: on the system's own price data, a name-based query for
``cement'' alone matches 135 products spanning 1,400 to 4,766,000 VND -- a
factor of more than 3,400 -- because unit and product grade are conflated.
``What does cement cost'' has no single answer; a RAG-only system would pick
the semantically closest passage and present it as if it were correct
(Section 3.2.2 gives further evidence).
\end{quote}

\begin{quote}
\textbf{Principle 2 -- ``Not found'' is always preferable to ``wrong.''}
When no row satisfies all conditions, a tool must return a not-found status
rather than the closest match; when RAG evidence is insufficient, the model
must refuse or ask for clarification. This principle drives several concrete
choices detailed in Sections 3.2--3.3: name-matching relaxation stops at two
words rather than one, tokens containing digits are never dropped during
relaxation, column labels are never emitted under uncertainty, an
unreconciled OCR numeric cell is blanked rather than kept, and unit prices
outside a plausible range are excluded before being used in a cost estimate.
\end{quote}

\subsubsection{1.3 Core and Integrated Components}

The project's research contribution centers on the document-ingestion,
structured-RAG, and pricing-tool paths, where the Chapter~VI benchmark
results directly inform design. The voice, web-search, and deep-research
modules are implemented and running in the product, and are described at a
level sufficient to present the full architecture, but they are not the
object of the quantitative benchmark in Chapter~VI.

\begin{table}[htbp]
\centering
\caption{Core versus integrated components}
\label{tab:core-vs-integrated}
\small
\renewcommand{\arraystretch}{1.3}
\begin{tabularx}{\textwidth}{|p{0.28\textwidth}|X|p{0.16\textwidth}|}
\hline
\rowcolor{tableheader}
\textbf{Group} & \textbf{Components} & \textbf{Section} \\
\hline
Core -- quantitatively benchmarked &
Dual-library PDF extraction (text + table geometry), table resolver &
Section 3.2.1 \\
\hline
Core -- quantitatively benchmarked &
Table-aware chunking, hybrid (dense + BM25) vector store &
Section 3.2.2 \\
\hline
Core -- quantitatively benchmarked &
Structured price table, tool-calling &
Section 3.2.3 \\
\hline
Operational infrastructure &
Streaming API, asynchronous ingestion queue &
Section 3.3 \\
\hline
Product integration &
Voice (STT/TTS), web search, deep research &
Section 4 \\
\hline
\end{tabularx}
\end{table}
\FloatBarrier

Figure~\ref{fig:rag-architecture} summarizes how the two data paths share a
single canonical extraction layer at ingestion time and how a request router
selects between them at query time; each element is developed in
Section~3.2.

\begin{figure}[htbp]
\centering
\begin{tikzpicture}[
  font=\scriptsize,
  box/.style={draw, rounded corners=2pt, minimum height=0.65cm, align=center,
    fill=orange!8, inner sep=4pt, text width=2.5cm},
  dbbox/.style={box, fill=blue!8},
  decision/.style={draw, diamond, aspect=2.4, align=center, fill=gray!10,
    inner sep=2pt, font=\scriptsize},
  arr/.style={-{Stealth[length=2mm]}, thick},
  lbl/.style={font=\scriptsize\itshape}
]

\node[box, fill=green!10] (pdf) {PDF price bulletin};
\node[box, right=0.5cm of pdf, text width=3.2cm] (dual) {Dual-library parsing:
text blocks (coords) + table geometry};
\node[box, right=0.5cm of dual, fill=yellow!15, text width=3cm] (canon)
  {Canonical table grid (shared source of truth)};

\node[box, below left=1.0cm and -0.7cm of canon, text width=3cm] (branchA)
  {\textbf{Branch A}: HTML chunk, header repeat, T1500};
\node[box, below right=1.0cm and -0.7cm of canon, text width=3cm] (branchB)
  {\textbf{Branch B}: row extraction, name/unit/price parsing};

\node[dbbox, below=0.9cm of branchA] (vecdb) {Vector store\\dense + BM25};
\node[dbbox, below=0.9cm of branchB] (pricedb) {Relational DB\\material\_prices};

\node[decision, below=1.5cm of canon, xshift=0cm] (router)
  {Request Router\\(runs \emph{before} retrieval)};

\node[box, below=1.6cm of router, fill=red!8, text width=3.6cm] (principle)
  {Backend calls the pricing tool directly for numeric questions --\\
  \textbf{no RAG context is built first}};

\draw[arr] (pdf) -- (dual);
\draw[arr] (dual) -- (canon);
\draw[arr] (canon) -- (branchA);
\draw[arr] (canon) -- (branchB);
\draw[arr] (branchA) -- (vecdb);
\draw[arr] (branchB) -- (pricedb);

\draw[arr] (vecdb.south) -- (router.north west);
\draw[arr] (pricedb.south) -- (router.north east);
\draw[arr] (router) -- (principle);

\node[lbl] at (canon.north -| pdf.west) [anchor=south west, yshift=0.5cm]
  {ingestion time};
\node[lbl] at (vecdb.south -| pdf.west) [anchor=north west, yshift=-0.1cm]
  {query time};

\end{tikzpicture}
\caption{Shared canonical extraction feeding two complementary data paths,
with query-time routing decided before any retrieval runs.}
\label{fig:rag-architecture}
\end{figure}
\FloatBarrier

\subsection{2. ASR Component Design}
\guidance{To be completed by the Speech AI Lead. Suggested structure:}

\subsubsection{2.1 Model Serving and Inference}
\guidance{Describe how PhoWhisper and the RouteLoRA adapters are loaded and
served (in-process on CPU vs. a dedicated GPU host), and any latency
optimizations.}

\subsubsection{2.2 Adapter Loading and Routing Design}
\guidance{Describe how regional adapters are stored, selected, or combined
at inference time, and how routing confidence is computed and logged.}

\subsubsection{2.3 Audio Preprocessing Pipeline}
\guidance{Describe audio capture, resampling, noise handling, and the
interface contract (input/output format) with the upstream voice channel and
the downstream Agentic RAG component.}

\subsection{3. Agentic RAG Component Design}

\subsubsection{3.1 AI Model Integration}

Rather than using a single model for every task, the system deliberately
assigns a different model to each responsibility: one embedding model
dedicated to producing vectors; one primary chat model that reads context and
performs tool calling; several small, inexpensive classifier models that
handle narrow-scope tasks such as topic guarding, follow-up-question
rewriting, and disambiguating ambiguous name-matching candidates; vision
models that participate only in the OCR fallback path; and all arithmetic
and structured-query execution running in application code rather than being
delegated to any model.

The model that answers customer questions in production is configured in
exactly \textbf{one} place and shared by every response path -- RAG
streaming, price-lookup presentation, cost-estimate presentation, and the
tool-calling loop in Agentic mode -- and the user interface reads back that
same configuration value, so there is no risk of the interface displaying
one model name while a different model actually runs. This separation exists
for three reasons: (1)~\emph{benchmark reproducibility} -- re-running
evaluation with a different model is a different experiment, not a
configuration update, so the model identifiers that appear in the Chapter~VI
result tables are kept as \emph{historical evaluation artefacts} rather than
updated to match the production model; (2)~\emph{operating cost and product
stability} -- the production model choice weighs per-request cost, latency,
and provider availability, factors that a controlled generation benchmark
(which holds retrieval fixed and only swaps the model) does not measure; and
(3)~\emph{consistency across the response pipeline} -- if every feature
picked its own model based on its own benchmark, the system would be
difficult to audit and to explain when something goes wrong. A direct
architectural consequence is that \textbf{the representation used for
indexing} (HTML, locked in by the Recall@5 evidence in Section~3.2.2) and
\textbf{the model used to generate answers} (an operational configuration
parameter, changeable without re-ingesting data or recomputing embeddings)
are two independent decisions, even though both originate from the same
experimental chapter.

Tool calling is the mechanism that lets the chat model act on structured
data it cannot access directly: available functions and their parameters are
declared to the model in advance; when the model determines a function is
needed, it returns a structured call (function name and arguments) instead
of free text; the backend executes the corresponding logic (typically a
database query or a calculation) and returns the result as a separate
message for the model to read and phrase. Critically, the model cannot
fabricate a number this way -- it only sees what the function actually
returned, and if the function reports no result, the model's only available
content is ``no data found.''

\subsubsection{3.2 Data Flow and Processing}

\paragraph{3.2.1 Document ingestion and canonical extraction.}
Uploading a document does not block on the full pipeline (table parsing,
chunking, embedding, and price extraction can take from several seconds to
several minutes); the upload request instead enqueues the job and returns
immediately, and a bounded, sequential queue keeps a single large document
from stalling the HTTP layer or monopolizing application resources.

Each knowledge base carries two independent configuration flags. The
\emph{price-extraction flag} controls whether an upload also extracts
structured price rows into the relational database; because every price
query filters by region, a row saved without a region could never be found,
so the region field is required at upload time and the system rejects the
request immediately rather than silently ingesting dead data. The
\emph{table-heavy chunking flag} selects which chunking profile a document
uses (Section~3.2.2); changing the flag does not retroactively re-chunk
already-ingested documents. Both processing branches must carry the same
complete chunking configuration in the queue message -- an earlier version
sent the price-extraction branch only region and period, omitting chunk
parameters entirely, so the one knowledge-base type that most needed a
specialized profile was the only type that could not configure one; this
incident illustrates a general principle for multi-branch asynchronous
pipelines: every branch that processes the same document must carry the same
complete configuration, not a subset judged ``sufficient'' for the branch
being written at the time.

PDF extraction deliberately uses two specialized libraries rather than a
single general-purpose one: a coordinate-based text extractor is fast at
pulling text blocks and their positions, while a table-geometry extractor is
what makes it possible to \textbf{distinguish a genuinely empty cell from
the body of a vertically merged cell} using each cell's geometric
information -- a capability with no substitute. Skipping it has a concrete
failure mode: a ``technical standard'' column merged across a twelve-product
family collapses onto only the first product, and the language model fills
the rest with generic background knowledge, producing wrong answers for the
other eleven products.

Because the text extractor does not know which regions are tables, it
returns all text -- including text inside tables -- as unordered blocks;
without correction, the same data would appear twice (once in an HTML table
chunk, once in a scrambled text chunk), and fragments of one row's numbers
could land in the surrounding context of an adjacent table chunk, causing a
price from one row to be misread as belonging to another row with an empty
price cell. The pipeline therefore processes content in true reading order:
table regions are located geometrically first, any text block whose center
falls inside a table region is removed from the text stream, and the
remaining table and text segments are reordered by page coordinates.
Measured across the ten source PDFs, this filtering reduced total chunk
count from over three thousand to about twelve hundred -- primarily removing
duplication -- while the table-chunk count was unchanged.

Constructing the \textbf{canonical table grid} is the technically hardest
step and is shared by both data paths. It must distinguish three cases that
look identical if only cell text is inspected: a cell with a real value; a
genuinely empty cell (bordered, no text); and a position inside the body of
a cell merged vertically across several rows above it. The distinguishing
signal comes from table geometry (borders, per-cell bounding boxes), not
from cell text. A further, more dangerous case is a malformed grid: some
PDFs have incomplete table borders, so a cell that does contain real text is
dropped entirely because the table-reading tool cannot construct a cell at
that position -- and geometrically this looks \emph{identical} to the body
of a merged cell. Applying the ordinary inheritance rule in that case would
fill the position with a value borrowed from an unrelated row: not missing
data, but \textbf{fabricated} data that looks well-formed (correct numeric
format, correct unit, only the value is wrong). Before this was caught and
fixed, 121 real data rows carried fabricated values of this kind. The fix
re-examines the page directly: it reconstructs the spatial box of the
suspect position from column and row boundaries and checks whether any word
on the page actually falls inside that box -- if so, the position holds
genuine data for that row and is not inherited; if not, it is genuinely the
body of a merged cell and inheritance proceeds normally. This is a direct
application of Principle~2: a geometric inference that is correct in the
majority of cases can still fabricate data in the remaining minority, and
the cost of eliminating that minority -- one additional text scan of the
page -- is worth paying against the risk of a wrong number entering the
price database.

\textbf{Header recovery for continuation pages} follows a priority order: if
the first row of a new page matches the currently held header, repeat it; if
the column count matches the previous table and the first row clearly has
the shape of data (no column-label keyword present), treat the page as a
continuation and borrow the previous header; if the first row shows genuine
header characteristics (short cells, at least one familiar column-label
keyword, no cell that is purely a currency amount), create a new header; and
if none of these applies with sufficient confidence, \textbf{no header
label is emitted} -- a missing label only degrades quality, while a wrong
label can silently change the meaning of an entire column for every row
beneath it, another direct application of Principle~2.

For table chunks, \textbf{surrounding context} (a token-budgeted span of
text immediately before and after the table, cut at sentence boundaries) is
attached to restore ``what table is this, under what pricing conditions''
for both embedding and generation, but only after table-region text has
already been filtered from the text stream (above); otherwise the
surrounding context itself becomes a scrambled sequence of table cells,
which is more dangerous than no context at all. This mechanism is
\textbf{disabled} for near-all-table price annexes, where neighbouring
content is typically only repeated page header/footer noise, and kept
\textbf{enabled} for mixed prose-and-table documents, where an explanatory
paragraph commonly sits beside the table it describes.

An oversized table chunk that exceeds the embedding API's hard token ceiling
is split further at row boundaries with the header repeated in every
resulting piece; a single row that alone exceeds the budget (for example, one
cell containing a long free-text note) is dropped rather than corrupting the
whole embedding batch. A \textbf{visual OCR fallback} activates only when a
document produces zero chunks through the CPU-first path -- a signal that the
PDF is a scanned image -- avoiding the cost of calling a vision model on
every page of documents that already carry a text layer. Each page is then
processed through \textbf{two independent passes using vision models from
two different providers}: one pass extracts the table as structured HTML,
the other independently re-reads the same page as plain text purely to
cross-check numbers; any amount that appears in the structured pass but is
\textbf{not} confirmed by the independent cross-check pass is blanked rather
than kept, because a blank cell honestly signals ``no data'' while a
fabricated number becomes a silent error downstream. Two different
providers are required because a single model asked to reproduce a table
tends to fabricate a plausible value for a missing cell, and asking the same
model to check itself typically reproduces the same fabricated value.

\paragraph{3.2.2 Branch A -- table-aware RAG for structural questions.}
Normalized tables are stored as HTML, kept as the canonical representation
because it preserves row--column structure, distinguishes headers from data,
is easy to audit, and can be deterministically rendered into other
representations without changing any cell value. The Chapter~VI
representation experiment gives this choice direct quantitative support: on
the same 2,252 logical T1500 chunks, HTML reaches Recall@5 = 343/500
(68.6\%), versus 306/500 (61.2\%) for key--value and 218/500 (43.6\%) for
verbalized prose -- so \textbf{HTML is the default representation for
embedding and indexing}, not merely a storage convenience. Key--value and
verbalized strings can still be rendered from the same canonical grid at
generation time if a deployed model's own benchmark shows a benefit, keeping
``representation good for retrieval'' and ``representation good for a given
model to read'' as separate concerns.

The target table-aware configuration is \textbf{T1500}: a nominal 1,500-token
limit, splitting only at row boundaries (a row is never split across two
chunks), the header repeated in every child chunk and counted against the
token budget, with source, page, chunk type, and business metadata stored
alongside. In deployment, T1500 and the T3000 reference configuration are
not two arbitrary options but map directly onto the two chunking profiles
selectable per knowledge base: a default \texttt{STANDARD} profile (3,000-token
ceiling, corresponding to T3000) suited to prose with small or moderate
embedded tables, where finer splitting was measured to leave each piece too
few rows and made same-table chunks harder to distinguish in embedding
space; and a specialized \texttt{TABLE\_HEAVY} profile (1,500-token ceiling,
corresponding to T1500, with surrounding-context attachment disabled) suited
to near-all-table price annexes, where there is no prose to borrow context
from. Note that a chunking profile does \emph{not} include a choice of
generation model -- the end-to-end configuration confirmed in Chapter~VI
also names a specific generation model, but that choice is a separate,
operational decision (Section~3.1), not part of the chunking profile itself.

Table chunks are embedded and stored in the vector store together with
document, knowledge-base, page, chunk-type, and business metadata (region,
publication period, source type). Dense retrieval alone is measurably weak
for product codes, proper names, sizes, and numeric strings -- two codes
differing by a single character can sit very close together in vector space
while their real-world values differ substantially -- so the structural
branch combines dense vector search over the HTML index with BM25~Okapi,
fused by Reciprocal Rank Fusion ($k=60$), each branch contributing 50
candidates before fusion, with the top~5 used as generation context. Adding
BM25 is the single largest measured improvement in the whole study: on the
500-question benchmark, Recall@5 for the table-aware branch rises from
68.6\% to 82.0\% -- a gain of 67 questions, larger than the entire gap
between table-aware and recursive chunking. Two safety gates govern the
fusion: metadata filters are applied to \emph{both} the dense and sparse
branches (otherwise BM25 could surface a chunk from the wrong region for a
region-qualified question, defeating a carefully designed region filter --
Section~3.4); and a similarity-score threshold is applied \emph{only} to the
dense branch, because after RRF fusion the combined score no longer shares a
scale with raw cosine similarity, so BM25 is restricted to re-ranking and
extending what dense retrieval already surfaced rather than answering on its
own when no dense candidate clears the threshold.

\paragraph{3.2.3 Branch B -- structured data and the pricing tool.}
A price lookup differs fundamentally from an attribute question: a price is
correct only when the product or code, specification, unit, region,
publication period, source type, and price basis all match simultaneously.
Vector search is optimized for semantic closeness, not exact matching on
these business fields, and it does not natively support aggregation (lowest,
highest, count, filter by period). These weaknesses are inherent to semantic
retrieval, not a matter of tuning: two product codes differing by one
character can be nearly identical in vector space while differing
substantially in price; a number and a material name can be misaligned
within the same chunk when a neighbouring row's leftover context is folded
in; and no aggregation operator exists in vector search at all -- it simply
returns $k$ passages and stops. Quantitative evidence on more than ten
thousand extracted price rows makes the case concretely: a name search for
``cement'' alone matches 135 products spanning 1,400 to 4,766,000~VND, and
a search for ``steel'' matches 512 products spanning 3 to
210,000,000~VND -- differences of several orders of magnitude from
conflating unit and grade in one descriptive query. Price is therefore
extracted into structured data and queried through a tool; RAG supplies only
supplementary explanation and citation, never the source of a number in the
final answer.

Price-row extraction runs on the same canonical grid built in
Section~3.2.1, so the RAG and SQL paths never see two different versions of
the same table. Header detection scans the first several rows for
keyword matches per column type (name, unit, group, price variants), merges
headers that span multiple physical rows through a scanning window, and
skips (with a logged warning) any table where at minimum a name column, a
unit column, and at least one price column cannot be identified -- rather
than guessing. When direct detection and borrowing the column mapping from
the previous page both fail, a fallback layer infers columns purely from
\textbf{data shape}, independent of labels: a column whose cells mostly
match a currency format is the price column, a column whose cells are mostly
a single word from a familiar unit vocabulary is the unit column, and the
longest text column preceding the unit column is the name column; missing
any of these three, the table is still skipped. Row-level extraction must
also handle group-header rows that masquerade as products once merged cells
are filled in, resolve repeated-value symbols in the unit column, prevent a
material-group label from bleeding into the next group, only accept a
manufacturer name when the value has the shape of an organization name,
normalize numbers to Vietnamese convention despite font kerning artifacts,
and split a single row into multiple price records when parallel price
columns exist (for example, ex-factory price and delivered price). A row
with a name and unit but no readable numeric value creates no record and is
pushed to a warning list instead of a guess, since a row with a wrong price
is more dangerous than a row that is simply missing.

Material-name matching is the largest bottleneck on the price path, measured
on a set of sixteen real lookup queries. The previous substring-matching
approach required the user's words to appear contiguous and in order, which
fails whenever the stored name is more complete than the user's shortened
query. The current approach: tokenize the query, strip diacritics and stop
words; require a candidate to contain \emph{all} remaining tokens (not any),
while allowing other tokens to appear interspersed -- as strict as before on
content, more flexible on order; rank surviving candidates by string
similarity; and if no row contains every token, progressively drop the most
common remaining token and retry. Measured first-result accuracy rose from
5/16 to 15/16; the one remaining miss is a material genuinely absent from
the queried region's data, where returning no result is correct. Two safety
gates were added only after a measurement caught a real failure, not as
theoretical caution: tokens containing digits are never dropped during
relaxation, since product codes and dimensions carry a material's identity;
and relaxation never proceeds below two remaining tokens, since a single
token is not sufficient identity.

The pricing tool exposes region, material name, and material group as
parameters, and deliberately marks \textbf{none of them as required} -- when
region was previously required, the model tended to guess a region for a
question that did not name one, directly violating Principle~2; missing-slot
follow-up is instead handled entirely by the query-time router
(Section~3.2.4). Typed, parameter-based tool-calling is used rather than
letting the model generate its own query (text-to-SQL) at the current stage,
for four verifiable reasons: the query space is currently small and fully
describable by a parameter schema; a typed call can be tested by fixing
parameters and comparing output, while a model-generated query differs every
run and is hard to reproduce for debugging; the security surface is
essentially absent, since the model never gains direct access to tables
holding user data or credentials; and, most importantly, the actual measured
bottleneck is name matching (above) -- a model-generated query would still
need the same substring matching already shown to be insufficient, so
text-to-SQL would not touch the real problem at this stage.

Because the Chapter~VI rejection benchmark shows chunking and retrieval do
not reliably substitute for business validation (McNemar $p=0.6636$ between
table-aware and recursive rejection rates), the pricing path runs an
explicit chain of checks before emitting any number: whether the product or
code truly matches, whether the unit is appropriate, whether the requested
period exists, whether the requested field is empty, and whether the row
carries a valid document source. A lookup result is classified into one of
four structured states, each handled distinctly rather than with a generic
answer: exactly one matching row yields a fact table the model may only
restate; multiple matching candidates are listed for the user to choose from
rather than averaged or auto-selected; a missing required parameter (region
or material name) triggers a targeted follow-up question; and no matching
row returns an explicit not-found statement -- never a price from a
different region, never a value from RAG, never a similar substitute chosen
automatically.

\paragraph{3.2.4 Query-time orchestration.}
Principle~1 requires that a number in an answer come only from the
structured price table or a code-level calculation, never from the language
model itself -- which in turn requires the decision ``does this question
need an exact number'' to be made \emph{before} any RAG context block is
built, not left to the model to notice after being handed a large context.
This is not an abstract design preference; it follows from a concrete
controlled observation made during development. The original design let the
model decide for itself whether to call the pricing tool after RAG context
had already been supplied, assuming the model could reliably recognize when
an exact number was needed. A controlled test showed that assumption false:
with RAG context attached directly to the question, the tool was called
\textbf{zero times} -- the model treated the supplied passages as the
answer source, failed to find the exact product in them, and concluded
``no data'' without ever attempting a tool call; disabling RAG context for
the identical question caused the tool to be called immediately. In other
words, the retrieval mechanism designed to \emph{support} the exact-lookup
path was in practice \emph{disabling} it. Moving the RAG context into a
separate system message with an explicit ``context does not replace the
tool'' instruction improved matters but still rested on a fragile premise: a
prompt instruction is a \emph{suggestion}, while execution order in code is
not. The final fix removes the dependence on the model's own initiative
entirely: a \textbf{request router runs before any retrieval and before the
main model is called}, deciding whether a question needs the deterministic
numeric path, the RAG path, or both; for the numeric path the backend calls
the pricing tool directly without asking the model's permission, and for
exact-price questions specifically, \textbf{no RAG retrieval runs at all}.

Early-exit checks answer cheaply before invoking the main model wherever
possible: small talk is matched \emph{exactly} (not by substring) against a
normalized set of greetings; a fixed-intent detector recognizes requests
that map onto an input form (typically a construction cost estimate) via
keyword-group matching; a topic guard (only active when no specific
knowledge base or project is selected) calls a small classifier to reject
off-topic questions politely, defaulting to pass-through on error rather
than blocking incorrectly; and a follow-up question is rewritten into a
self-contained query using recent conversation history before routing or
retrieval. The router itself returns a structured decision -- processing
route, intent, named regions, price period, material name and group,
manufacturer, requested fields, missing slots, and a confidence score --
governed by one architectural rule: \textbf{deterministic keyword rules run
before any classifier, and once a deterministic rule has identified a clear
price question, a classifier is never permitted to downgrade it to a plain
RAG question.} A model classifier is invoked only as a fallback when
deterministic rules cannot decide, and on classifier error the system
defaults to ordinary conversation rather than hard-blocking. Table~\ref{tab:router-rules}
summarizes the main routes.

\begin{table}[htbp]
\centering
\caption{Query-time router: representative rules}
\label{tab:router-rules}
\small
\renewcommand{\arraystretch}{1.25}
\begin{tabularx}{\textwidth}{|X|p{0.28\textwidth}|}
\hline
\rowcolor{tableheader}
\textbf{Question pattern} & \textbf{Route} \\
\hline
Asks the price or unit price of a specific material & \texttt{EXACT\_STRUCTURED} \\
\hline
Asks the unit, manufacturer, specification, price basis, or period of a
specific product & \texttt{EXACT\_STRUCTURED} \\
\hline
Asks the material quantity or cost for a given floor area / project & \texttt{ESTIMATE} \\
\hline
Asks about VAT, scope of application, standards, or a comparison between
two different products & \texttt{DOCUMENT\_RAG} \\
\hline
Asks for both a price and a condition/standard in the same question & \texttt{MIXED} \\
\hline
Asks for a price without naming a region, or without naming a specific
material & \texttt{CLARIFY} \\
\hline
\end{tabularx}
\end{table}
\FloatBarrier

The user interface offers four interaction modes, each mapped to a different
backend retrieval scope and tool-access policy: \textbf{Chat} (default;
retrieval scoped to the selected knowledge base or project, with tool access
whenever the router determines a price/estimate question); \textbf{Search}
(web only, no internal tools); \textbf{Deep Research} (web, through a
multi-step loop, no internal tools); and \textbf{Agentic} (retrieval across
every knowledge base the user can access, with three public tools -- exact
price lookups are still forced through the tool path by the backend exactly
as in every other mode). The tool-calling loop in Agentic mode is bounded in
the number of turns and is deliberately \emph{not} given a free-form RAG
retrieval tool, since the backend orchestration layer already manages
retrieval and region filtering; granting a redundant retrieval tool would
only duplicate that work and risk bypassing the carefully designed region
filter (Section~3.4).

For an exact price question, the flow is: router determines
\texttt{EXACT\_STRUCTURED} $\rightarrow$ backend calls the pricing service
directly (no RAG retrieval) $\rightarrow$ service matches the name and
resolves candidates $\rightarrow$ a found row becomes a fact table that the
model may only restate at low temperature, or a not-found status is
returned verbatim without invoking the model at all. For a mixed question,
the tool and retrieval paths run \emph{in parallel} but do not carry equal
authority: the tool supplies price, unit, period, region, and provenance,
while RAG supplies VAT notes and scope-of-application text; structured
fields from the tool always take precedence, and RAG is never permitted to
overwrite a price, region, or unit produced by the tool. Beyond backend
mechanisms, the chat model's system prompt carries a final layer of guard
rules: never fabricate a price without accompanying real data; match the
exact requested region and say plainly when the requested region lacks data
rather than substituting another region's price; prefer supplied context
over background knowledge; decline questions outside the construction-
material domain even when a keyword happens to overlap; and, for a vague
follow-up question against narrow underlying data, ask for clarification
rather than mechanically listing whatever exists.

\subsubsection{3.3 Deployment Strategy}

The CPU-first document-extraction path (Section~3.2.1) runs entirely without
a vision model for text-layer PDFs, preserving per-cell coordinates and
geometry, reducing ingestion cost and latency, and remaining easy to audit
because every step can be re-checked against the original PDF page. Stronger
visual document-understanding models could help with scanned documents or
very complex layouts, but they substantially increase image size, memory,
latency, and operating cost for \emph{every} document, including the
majority that already carry a text layer and do not need that capability;
OCR therefore runs only as a conditionally triggered fallback
(Section~3.2.1), never as the default path. The specific pair of vision
models used for the OCR fallback's structured pass and cross-check pass was
itself chosen by direct measurement against a real scanned document from the
corpus, scored line-by-line against the printed original: the chosen
configuration matched the accuracy of a substantially more expensive model
at roughly an order of magnitude lower per-page cost.

Two storage systems are used because they serve non-overlapping roles,
summarized in Table~\ref{tab:storage-roles}. The relational database is the
system's ``source of truth for numbers''; the vector store is its ``semantic
memory and evidence store.'' A vector store (rather than a vector extension
inside the existing relational database) was chosen for three
problem-specific reasons: metadata filtering (knowledge base, region) within
a single query rather than filtering after retrieval; native support for
both dense and sparse search to enable the hybrid retrieval of
Section~3.2.2; and the ability to operate independently without locking the
system to one cloud provider. The in-database vector-extension alternative
was considered -- simpler operationally, since it avoids running a
separate service -- but metadata filtering is more complex and no native
sparse-retrieval mechanism is available.

\begin{table}[htbp]
\centering
\caption{Division of responsibility between the two storage systems}
\label{tab:storage-roles}
\small
\renewcommand{\arraystretch}{1.3}
\begin{tabularx}{\textwidth}{|p{0.2\textwidth}|X|X|p{0.16\textwidth}|}
\hline
\rowcolor{tableheader}
\textbf{Store} & \textbf{Data} & \textbf{Query type} & \textbf{Expected precision} \\
\hline
Relational database & Normalized price rows & Filter, sort, aggregate,
deterministic matching & Deterministic \\
\hline
Vector store & Text and table chunks & Semantic search + metadata filter &
Approximate \\
\hline
\end{tabularx}
\end{table}
\FloatBarrier

An asynchronous queue is used for document ingestion because a long
document is a multi-step task that can run for minutes: it gives a fast API
response (returns a job ID immediately instead of waiting), automatic retry
on worker failure, a bound on concurrent jobs to avoid resource contention,
and it decouples the lifetime of an HTTP request from the lifetime of the
full ingestion process, which is far longer than a connection should
reasonably be held open. A heavier dedicated task-queue system (with a
separate broker and result store) was judged more than the need warranted;
running background work directly inside the web-application process was
rejected because in-flight jobs would be lost on every process restart --
unacceptable for work that can run for several minutes.

Technology choices for the integrated components (Section~4) follow the
same philosophy: Vietnamese speech recognition uses a model fine-tuned
specifically for Vietnamese, capable of running on CPU in-process (with an
optional dedicated GPU host for lower latency), avoiding dependence on a
usage-billed cloud speech service; speech synthesis calls a dedicated
external service directly, since this is not a core capability that needs
to be self-hosted; and web search / deep research use a dedicated
search-and-scrape service together with a graph-structured multi-step
orchestration framework rather than a custom-built crawler and loop.

\subsubsection{3.4 Scalability and Maintenance}

\textbf{Configuration identity and cache.} An earlier configuration-alias
mix-up (Section~4, Chapter~IV) motivated attaching an explicit identity to
every cached or logged result at both the research and production layers:
an explicit configuration name for each chunking profile, a configuration
``signature'' composed of chunk count and content digest to confirm a run or
cache retrieval is using the expected data, a rule that a vector cache is
reused only on a full digest match, separate storage paths for raw and
aggregated results, and a generation record that carries the query
identifier, the chunk list actually used, a prompt hash, model name, and
stop reason. These safeguards are part of the measurement architecture
itself, not benchmark-only detail -- the same ``attach configuration identity
to every result'' principle also governs production log tracking below.

\textbf{Citation and provenance.} Every answer must trace back to the
correct source -- one of the five design goals in Section~1.2. To do this
consistently across the tool and RAG paths, every citation is normalized to
a common structure regardless of origin (a tool lookup, a RAG chunk, or a
web search result): source type (tool/RAG/web), authority level
(verified/supporting/unverified), purpose (price/structured field/
explanation/estimate), plus document, page, chunk or row identifiers,
region, publication period, and relevance score.

A notable production incident illustrates why source normalization matters
more than it appears to. Symptom: a user asked for a price in a specific
region; the response text could be correct, but the source badge shown on
the interface named a completely different region -- even though the
underlying data had been correctly tagged with region at ingestion time.
Tracing the region field through every hop -- upload, document metadata,
chunk metadata, vector-store payload, out through the retrieval layer --
showed the region field \textbf{survived intact at every hop except exactly
one}: the step that packages the source list for the interface, which
serialized only file name, content, and score, silently dropping the region
field that was available one step earlier. Because the interface had
nothing else to show, a file name that happened to contain the name of a
different region was misread as the answer's region label. More
significantly, this was \textbf{not simply a labeling bug}: the retrieved
chunk itself genuinely belonged to the wrong region, because region
filtering at retrieval time was, at that point, enabled only conditionally
for a specific knowledge base -- so in Agentic mode (retrieval across every
accessible knowledge base) and project-scoped chat (multiple knowledge bases
merged), \textbf{no region filter ran at all}, meaning the mislabeled badge
was only the visible symptom of a deeper gap through which wrong-region data
could genuinely reach the answer context in those two modes. The fix
required two independent corrections: ensuring region information survives
every serialization step without exception, and moving the region-filter
condition away from ``which knowledge base is currently selected'' toward
``which region is named in the request itself,'' so filtering behaves
consistently regardless of interaction mode. Source rules distilled directly
from this incident now apply to every response path: never infer region
from a file name, knowledge-base name, answer content, or the user's
previous stated region; a region-neutral source may be kept only when it is
genuinely neutral, never relabeled to look correct; multi-region comparison
questions retrieve separately per region and include region in the
deduplication key; and the interface badge (e.g., ``Verified lookup'' versus
``RAG'') is decided by which source types are genuinely present, not merely
by whether a source field exists at all.

\textbf{Observability and data cleanup.} Each chat turn logs a structured
record for routing and region-filter debugging (chosen route, tool status,
source counts before/after region filtering), without access credentials or
sensitive content; the system additionally tracks ingestion time and chunk
count per document, extracted price-row counts, extraction warning counts,
token usage and model cost, retrieval and generation latency, tool
not-found rate, and refusal rate. Because the relational database and the
vector store are two independent systems with no shared transaction
(Section~3.3), deleting a knowledge base or a document must actively clean
both sides in a specific order -- verify ownership, delete vectors first,
delete relational data second -- because the two possible half-completed
failure states are not equally severe: losing vectors while relational data
survives is a \emph{visible}, recoverable error (re-ingest); losing
relational data while vectors survive is an \emph{invisible} error, since
the system continues to answer normally while silently relying on data the
user believes has been deleted. When forced to choose between failure
modes, the architecture chooses the one that can be recovered.

Selected production incidents, recorded deliberately to avoid repeating them
when designing new features, are summarized in
Table~\ref{tab:production-incidents}.

\begin{table}[htbp]
\centering
\caption{Selected production incidents and the resulting architectural lesson}
\label{tab:production-incidents}
\footnotesize
\renewcommand{\arraystretch}{1.3}
\begin{tabularx}{\textwidth}{|p{0.16\textwidth}|X|X|}
\hline
\rowcolor{tableheader}
\textbf{Incident} & \textbf{Root cause} & \textbf{Architectural lesson} \\
\hline
Orphaned vectors after knowledge-base deletion &
Relational cascade delete has no knowledge of the vector store &
Two independent stores require explicit application-level cleanup; cascade
in one store never spans the other (Section~3.3) \\
\hline
Source badge showing the wrong region &
Region field dropped at one serialization step; region filtering gated on
selected knowledge base rather than the request's own region &
Normalize the source schema end to end and filter by the request's region,
not the selected scope (Section~3.4) \\
\hline
Pricing tool never called &
RAG context attached directly to the question caused the model to treat it
as the only answer source &
The decision to call a tool must sit in the backend and run before
retrieval, never left to the model to notice (Section~3.2.4) \\
\hline
Fabricated price from a malformed table grid &
A position with no cell due to missing border lines was indistinguishable,
by geometry alone, from a merged-cell body &
Two failure modes that share one surface signal need an independent
verification step, not a single-signal inference (Section~3.2.1) \\
\hline
Unrealistic unit price entering a cost estimate &
No plausibility bound was applied to a web-sourced unit price &
Any lower-confidence numeric source, especially real-time web data, needs a
plausible-range check before entering a calculation (Section~4.3) \\
\hline
\end{tabularx}
\end{table}
\FloatBarrier

\textbf{Tool security.} The typed tool-calling mechanism (Section~3.2.3)
doubles as an access-surface control: the model is never granted a direct
database connection or free-form SQL execution; every tool-backed function
performs only its predefined queries against exactly the tables it needs,
with bounded result sizes, following the no-guessing rule of
Section~3.2.3. This is a security benefit that follows from choosing
typed tool-calling over text-to-SQL, not a separately engineered security
layer.

\subsection{4. End-to-End Integration}

The system declares six tools in a shared tool server, but exposes only
three of them to the chat model in ordinary conversation, summarized in
Table~\ref{tab:tool-exposure}.

\begin{table}[htbp]
\centering
\caption{Declared tools and chat-agent exposure}
\label{tab:tool-exposure}
\small
\renewcommand{\arraystretch}{1.25}
\begin{tabularx}{\textwidth}{|X|>{\centering\arraybackslash}p{0.14\textwidth}|X|}
\hline
\rowcolor{tableheader}
\textbf{Tool} & \textbf{Exposed to chat agent?} & \textbf{Touches} \\
\hline
Material unit-price lookup & Yes & Relational database \\
\hline
Quantity estimation (takeoff formula) & Yes & Calculation only \\
\hline
Construction cost estimate from floor area & Yes & Relational database + web
(conditional) \\
\hline
RAG retrieval within a knowledge base & No & Vector store \\
\hline
Web search & No & External search service \\
\hline
Deep research & No & Multi-step orchestration + web \\
\hline
\end{tabularx}
\end{table}
\FloatBarrier

The three tools not exposed to the chat agent remain fully implemented and
reachable through a separate standardized protocol for external client
applications: a retrieval tool is redundant for the chat agent because
backend orchestration already manages retrieval and region filtering
(Section~3.2.4); web search and deep research have their own dedicated
interaction channels where the user actively opts in (below), so letting the
model call them autonomously mid-conversation could turn a
few-second response into one taking tens of seconds without warning.

The system supports two-way voice interaction. \guidance{This is the
integration boundary with the ASR component: the Speech AI Lead should
document the exact audio input format, expected latency budget, and the
transcript/confidence-score contract handed to the Agentic RAG orchestrator
described above, plus how a low-confidence transcription is surfaced (e.g.,
routed to \texttt{CLARIFY} rather than answered directly).} On the
production side, when a question arrives as speech, the system skips the
fixed-intent-detection step (Section~3.2.4) entirely, so a spoken utterance
always proceeds directly into the ordinary RAG/routing path rather than
risking an abrupt input form appearing mid-conversation, where filling out a
form is a poor fit for a voice interaction. Speech synthesis reads back
exactly the text the model generated, not separately authored
voice-specific content.

The \textbf{construction cost estimation tool} solves a different problem
from plain price lookup: an order-of-magnitude estimate of \emph{primary
material} cost for a project based on floor area, at a stage before design
drawings exist -- it excludes labour, equipment, contractor margin, and VAT,
and its precision matches a conceptual estimate rather than a detailed
quantity takeoff. Material quantities follow
$\text{quantity}_i = \text{reference area} \times \text{consumption factor}_i
\times \text{finish factor (if applicable)}$, with cost
$= \text{quantity} \times \text{unit price looked up from the price table}$
and total cost summed across materials; consumption factors and the
material set differ by building type (reinforced-concrete housing,
prefabricated steel structures, concrete yards, ground leveling, and so on),
since one shared coefficient set cannot serve both a townhouse and a
prefabricated steel warehouse. For each material, the system queries
candidate prices concurrently by region, name, and unit, and a small
classifier model selects exactly one correct row among the candidates (or
confirms none is suitable) -- a decision grounded in real database rows, not
a model-generated price. Missing a database price, the system
\textbf{fails closed by default}: it lists the missing line items and
withholds a total rather than silently omitting the item from an
understated sum; only with explicit user permission does it look up a
reference price from the web, always labeled clearly as unverified and never
silently blended with an officially published price. Two additional safety
gates, both added after real incidents rather than as theoretical caution,
apply to unit prices from either source: a \emph{plausible price band} per
material excludes out-of-range candidates (added after a web-sourced steel
unit price -- actually a whole-project total misread as a per-unit price --
inflated one line item's cost by several orders of magnitude while still
displaying with the same confident appearance as correct rows); and
\emph{alternate-unit conversion}, since the same material can be published
under more than one unit (e.g., by volume and by weight) -- querying only
one fixed unit would make a meaningful share of valid data invisible, so the
system queries the primary unit first, then an alternate unit with a
conversion factor if the primary unit is not found. The final presentation
gives a cost \emph{range} rather than a single figure when enough major line
items have prices, states clearly what is missing when at least one major
item is unpriced rather than presenting a partial total, and can invert the
calculation to suggest a feasible floor area from a user-supplied target
budget; the model is only asked to restate the already-fixed result, never
to recompute or adjust any figure.

Web search and deep research serve information needs outside the ingested
document set and have their own interaction channels on the interface --
selected deliberately by the user rather than triggered autonomously by the
model mid-conversation (Section~3.2.4). Deep research is a bounded
multi-step loop -- rewrite into an independent query, expand into search
directions, fetch and collect web content, synthesize, check a quality
threshold, and either loop again or generate a cited answer -- stopping at a
quality threshold or a maximum number of iterations, whichever comes first,
so that response time stays predictable even for difficult questions, at
the cost of occasionally falling short of an ideal synthesis quality on the
hardest ones.

% -----------------------------------------------------------------------------
% VI. Results and Discussion
% -----------------------------------------------------------------------------
\chapterstart{VI. Results and Discussion}

This chapter reports the outcome of the evaluation methodology defined in
Chapter~IV. The Agentic RAG results reflect a completed benchmark on the
500-question and 70-question sets described in Section~IV.3.2; the ASR
results are left as a structured template for the Speech AI Lead to
complete once the corresponding experiments are run.

\subsection{1. Results and Analysis}

\subsubsection{1.1 ASR Recognition Results}
\guidance{To be completed by the Speech AI Lead. Suggested structure:
overall WER/CER for the PhoWhisper baseline versus RouteLoRA, broken down by
Northern/Central/Southern accent and by noise/domain-terminology condition;
a routing-confidence analysis; and a comparison against single-accent
adapters under known- and unknown-accent test conditions.}

\subsubsection{1.2 Agentic RAG Results}

All results below come from the final run with index identity locked by
chunk count, content digest, and cache version (Section~IV.3.5); all
generation was executed directly through a single inference endpoint, with
the full denominator of 500 questions retained throughout.

\paragraph{1.2.1 Structure preservation (RQ1).}
Table~\ref{tab:res-structure} compares structural quality across the three
main chunking configurations.

\begin{table}[htbp]
\centering
\caption{Structural quality of the main configurations}
\label{tab:res-structure}
\small
\renewcommand{\arraystretch}{1.25}
\begin{tabularx}{\textwidth}{|X|>{\centering\arraybackslash}p{0.11\textwidth}|>{\centering\arraybackslash}p{0.14\textwidth}|>{\centering\arraybackslash}p{0.12\textwidth}|>{\centering\arraybackslash}p{0.14\textwidth}|}
\hline
\rowcolor{tableheader}
\textbf{Configuration} & \textbf{Chunks} & \textbf{Avg.\ tokens} &
\textbf{Row integrity} & \textbf{Rows with column labels} \\
\hline
T1500 & 2,252 & 831 & 99.5\% & 100\% \\
\hline
T3000 & 1,476 & 1,185 & 99.5\% & 100\% \\
\hline
R1500 & 790 & 1,102 & 99.4\% & 65\% \\
\hline
\end{tabularx}
\end{table}
\FloatBarrier

After correcting for labels that did not appear evenly across branches,
T1500 and R1500 are nearly equivalent at keeping a material name and its
price in the same chunk (99.5\% versus 99.4\%); row integrity is therefore
\emph{not} used as the main argument for table-aware chunking. The clear
difference is in column labeling: T1500 and T3000 both retain the header in
100\% of table chunks, while R1500 carries sufficient header signal in only
about 65\% of rows under the current measurement. Recursive chunking can
keep values physically close together but no longer marks which cell
belongs to which column. This directly answers RQ1: the table-aware
pipeline's structural advantage is not superior row integrity, but retained
column labels and row--column relationships that allow the model to
interpret attributes correctly.

\paragraph{1.2.2 Cost in embedding space (RQ2).}
Table~\ref{tab:res-embedding} reports inter-chunk cosine similarity.

\begin{table}[htbp]
\centering
\caption{Inter-chunk cosine similarity by configuration}
\label{tab:res-embedding}
\small
\renewcommand{\arraystretch}{1.25}
\begin{tabularx}{\textwidth}{|X|>{\centering\arraybackslash}p{0.16\textwidth}|>{\centering\arraybackslash}p{0.12\textwidth}|>{\centering\arraybackslash}p{0.14\textwidth}|>{\centering\arraybackslash}p{0.14\textwidth}|}
\hline
\rowcolor{tableheader}
\textbf{Configuration} & \textbf{Median cosine} & \textbf{P90} &
\textbf{Share $\geq$0.90} & \textbf{Share $\geq$0.95} \\
\hline
T1500 & 0.633 & 0.788 & 1.67\% & 0.38\% \\
\hline
T3000 & 0.605 & 0.782 & 2.69\% & 0.65\% \\
\hline
R1500 & 0.611 & 0.731 & 0.59\% & 0.20\% \\
\hline
\end{tabularx}
\end{table}
\FloatBarrier

T1500 shows a share of pairs with cosine similarity $\geq0.90$ roughly 2.8
times higher than R1500 (1.67\% versus 0.59\%). Retaining and repeating
headers creates a shared span of text across chunks of the same table; this
shared text helps a model understand columns, but makes chunks harder for a
dense retriever to discriminate. T3000 shows an even higher share of
near-duplicate pairs than T1500 (2.69\% versus 1.67\%); larger chunks may
contain more rows with repeated vocabulary and structure, though this study
did not run an ablation separating header from body contribution, so this
explanation remains a post-hoc hypothesis. This answers RQ2: table-aware
chunking creates a genuine trade-off -- better preservation of headers and
column relationships, at the cost of an index with more near-duplicate
vector pairs than recursive chunking.

\paragraph{1.2.3 Retrieval profile at multiple cutoffs (RQ3, part 1).}
Extending measurement from Recall@5 alone to Recall@1/3/5/10 and MRR shows
that T1500 and R1500 have different ranking behavior; no single branch wins
consistently at every $k$ (Table~\ref{tab:res-retrieval}).

\begin{table}[htbp]
\centering
\caption{Retrieval at multiple cutoffs, full 500-question set}
\label{tab:res-retrieval}
\small
\renewcommand{\arraystretch}{1.25}
\begin{tabularx}{\textwidth}{|X|>{\centering\arraybackslash}p{0.1\textwidth}|>{\centering\arraybackslash}p{0.1\textwidth}|>{\centering\arraybackslash}p{0.1\textwidth}|>{\centering\arraybackslash}p{0.1\textwidth}|>{\centering\arraybackslash}p{0.09\textwidth}|>{\centering\arraybackslash}p{0.09\textwidth}|}
\hline
\rowcolor{tableheader}
\textbf{Config.} & \textbf{R@1} & \textbf{R@3} & \textbf{R@5} & \textbf{R@10} &
\textbf{MRR} & \textbf{No-hit@10} \\
\hline
T1500 dense & 41.4\% & 61.4\% & 68.6\% & \textbf{78.8\%} & 0.5345 & 106 \\
\hline
R1500 dense & \textbf{45.6\%} & \textbf{65.4\%} & \textbf{70.8\%} & 77.8\% &
\textbf{0.5659} & 111 \\
\hline
T1500 hybrid & 54.4\% & 74.6\% & \textbf{82.0\%} & 87.0\% & 0.6597 & 65 \\
\hline
R1500 hybrid & \textbf{58.8\%} & \textbf{76.2\%} & 80.6\% &
\textbf{87.6\%} & \textbf{0.6844} & 62 \\
\hline
\end{tabularx}
\end{table}
\FloatBarrier

In dense retrieval, R1500 is higher at R@1, R@3, R@5, and MRR; T1500 only
edges ahead slightly at R@10. This reinforces the conclusion that table-aware
chunking does \emph{not} create a general dense-retrieval advantage. In
hybrid retrieval, R1500 remains higher at R@1, R@3, R@10, and MRR, while
T1500 is higher at R@5 -- precisely the operating point the system uses to
feed context into generation. The precise statement is therefore that
\textbf{T1500 achieves a higher hybrid Recall@5}, not that ``T1500 has a
better retrieval ranking at every depth.'' Splitting Recall@5 by question
group makes the pattern sharper still: for price-lookup questions, hybrid
retrieval lifts T1500 from 189/252 to 238/252 and R1500 from 196/252 to
228/252, since identifier tokens (codes, names, sizes, numeric strings) suit
lexical matching; for attribute questions, R1500 remains higher than T1500
at Recall@5 in both dense (158/248 versus 154/248) and hybrid (175/248
versus 172/248) retrieval. So the attribute-generation advantage of the
table-aware pipeline (Section~1.2.6) cannot be explained by retrieval
finding evidence more often.

\paragraph{1.2.4 Effect of representation on dense retrieval (RQ3, part 2).}
Holding the 2,252 logical T1500 chunks fixed and varying only the string
passed to the embedding model isolates the effect of representation from
the effect of chunking (Table~\ref{tab:res-representation}); this experiment
does not call a generation model.

\begin{table}[htbp]
\centering
\caption{Recall@5 when the embedded representation changes, fixed T1500 chunks}
\label{tab:res-representation}
\small
\renewcommand{\arraystretch}{1.25}
\begin{tabularx}{\textwidth}{|X|>{\centering\arraybackslash}p{0.16\textwidth}|>{\centering\arraybackslash}p{0.12\textwidth}|>{\centering\arraybackslash}p{0.18\textwidth}|}
\hline
\rowcolor{tableheader}
\textbf{Representation} & \textbf{Recall@5} & \textbf{Rate} &
\textbf{vs.\ HTML} \\
\hline
\textbf{HTML} & \textbf{343/500} & \textbf{68.6\%} & -- \\
\hline
Key--value (KV) & 306/500 & 61.2\% & $-$37 questions ($-$7.4 pp) \\
\hline
Verbalized (VERB) & 218/500 & 43.6\% & $-$125 questions ($-$25.0 pp) \\
\hline
\end{tabularx}
\end{table}
\FloatBarrier

HTML achieves the highest observed Recall@5. Linearizing the same table data
into key--value form loses 37 retrieval hits; verbalized prose loses far
more, 125 hits. This result rejects the assumption that converting a table
to a ``more readable'' string necessarily improves embedding retrieval. It
also creates an important distinction from the generation benchmark in
Section~1.2.9: \emph{the representation best suited for indexing is not
necessarily the representation best suited for a given model to read}. At
retrieval, HTML is a clear winner at the R@5 endpoint; at generation, HTML
achieves the highest overall EM with GPT-OSS~20B and two Gemma models, while
KV achieves the highest EM with Llama~3.1~8B. The architecture therefore
still keeps \emph{retrieval representation} and \emph{generation
representation} conceptually separate, even though HTML is a reasonable
default for most models in the current benchmark.

\paragraph{1.2.5 Effect of BM25 on retrieval (RQ3, part 3).}
Table~\ref{tab:res-bm25} isolates the effect of adding BM25 at the top-5
operating point.

\begin{table}[htbp]
\centering
\caption{Retrieval $2\times2$ at the top-5 operating point}
\label{tab:res-bm25}
\small
\renewcommand{\arraystretch}{1.25}
\begin{tabularx}{\textwidth}{|X|>{\centering\arraybackslash}p{0.11\textwidth}|>{\centering\arraybackslash}p{0.14\textwidth}|>{\centering\arraybackslash}p{0.09\textwidth}|>{\centering\arraybackslash}p{0.13\textwidth}|>{\centering\arraybackslash}p{0.13\textwidth}|}
\hline
\rowcolor{tableheader}
\textbf{Config.} & \textbf{Dense R@5} & \textbf{Hybrid R@5} &
\textbf{Gain} & \textbf{Price-lookup gain} & \textbf{Attribute gain} \\
\hline
T1500 & 343/500 & 410/500 = 82.0\% & +67 & +49 & +18 \\
\hline
R1500 & 354/500 & 403/500 = 80.6\% & +49 & +32 & +17 \\
\hline
\end{tabularx}
\end{table}
\FloatBarrier

BM25 improves Recall@5 for both pipelines; T1500 gains 67 hits and R1500
gains 49. The difference-of-differences is $+18$ questions; a paired
bootstrap gives a 95\% confidence interval of $-2$ to $+37$ questions, with
$p=0.0796$ -- the interval contains zero, so there is not yet sufficient
evidence at the 0.05 threshold that BM25 interacts specially with
table-aware chunking. The gain is concentrated in price-lookup questions
(T1500: 189$\rightarrow$238, +49/252; R1500: 196$\rightarrow$228, +32/252),
with a smaller gain on attribute questions (T1500: 154$\rightarrow$172,
+18/248; R1500: 158$\rightarrow$175, +17/248). Combined with the ranking
profile in Section~1.2.3, BM25 does not make T1500 win at every cutoff --
R1500 hybrid still has higher R@1, R@3, R@10, and MRR. T1500 hybrid's value
lies specifically at the top-5 context the system uses and in the
attribute-generation results that follow.

\paragraph{1.2.6 Generation results (RQ4).}
Table~\ref{tab:res-generation} reports Exact Match for the four main
configurations.

\begin{table}[htbp]
\centering
\caption{Exact Match, four main configurations}
\label{tab:res-generation}
\small
\renewcommand{\arraystretch}{1.25}
\begin{tabularx}{\textwidth}{|X|>{\centering\arraybackslash}p{0.16\textwidth}|>{\centering\arraybackslash}p{0.18\textwidth}|>{\centering\arraybackslash}p{0.18\textwidth}|}
\hline
\rowcolor{tableheader}
\textbf{Config.} & \textbf{EM (overall)} & \textbf{Price-lookup G1--G4} &
\textbf{Attribute S1--S4} \\
\hline
T1500 dense & 251/500 = 50.2\% & 159/252 = 63.1\% & \textbf{92/248 = 37.1\%} \\
\hline
R1500 dense & 227/500 = 45.4\% & 162/252 = 64.3\% & 65/248 = 26.2\% \\
\hline
T1500 hybrid & \textbf{292/500 = 58.4\%} & \textbf{196/252 = 77.8\%} &
\textbf{96/248 = 38.7\%} \\
\hline
R1500 hybrid & 256/500 = 51.2\% & 182/252 = 72.2\% & 74/248 = 29.8\% \\
\hline
\end{tabularx}
\end{table}
\FloatBarrier

In dense retrieval, T1500 exceeds R1500 by 24 questions overall, entirely
driven by attribute questions (T1500 trails by 3 on price-lookup but leads
by 27 on attributes). In hybrid retrieval, T1500 exceeds R1500 by 36
questions overall (+14 price-lookup, +22 attribute). BM25 raises T1500's EM
by 41 questions (251$\rightarrow$292) and R1500's by 29 (227$\rightarrow$256);
for T1500, 37 of those 41 gained questions are price-lookup and only 4 are
attribute, while for R1500 the corresponding split is 20 price-lookup and 9
attribute. The observed generation interaction is $+12$ questions (T1500's
BM25 gain exceeds R1500's by 12), but a 100,000-resample paired bootstrap
gives a 95\% confidence interval of $-13$ to $+37$ questions with
$p=0.3723$ -- the interval contains zero, so there is not yet evidence that
BM25's effect on generation depends on the chunking strategy. Table-aware
chunking and BM25 resolve two different failure modes: table-aware chunking
creates its clearest advantage on attribute questions, where the model must
identify the correct column, while BM25 creates most of its benefit on
price-lookup questions, where product codes, sizes, proper names, and
numeric strings are highly discriminative.

\paragraph{1.2.7 Paired comparison and Holm correction.}
Table~\ref{tab:res-mcnemar} reports McNemar and Holm-adjusted p-values for
the primary confirmatory family of six T1500--R1500 comparisons.

\begin{table}[htbp]
\centering
\caption{McNemar and Holm-adjusted p-values, six-test family}
\label{tab:res-mcnemar}
\footnotesize
\renewcommand{\arraystretch}{1.25}
\begin{tabularx}{\textwidth}{|X|X|>{\centering\arraybackslash}p{0.06\textwidth}|>{\centering\arraybackslash}p{0.1\textwidth}|>{\centering\arraybackslash}p{0.1\textwidth}|>{\centering\arraybackslash}p{0.09\textwidth}|>{\centering\arraybackslash}p{0.09\textwidth}|}
\hline
\rowcolor{tableheader}
\textbf{Retriever} & \textbf{Scope} & \textbf{$n$} & \textbf{T$\checkmark$/R$\times$} &
\textbf{T$\times$/R$\checkmark$} & \textbf{$p$ raw} & \textbf{$p$ Holm} \\
\hline
Dense & Full set & 500 & 80 & 56 & 0.04818 & 0.14455 \\
\hline
Dense & Price-lookup & 252 & 36 & 39 & 0.8176 & 0.8176 \\
\hline
Dense & Attribute & 248 & 44 & 17 & \textbf{0.00073} & \textbf{0.00438} \\
\hline
Hybrid & Full set & 500 & 83 & 47 & \textbf{0.00202} & \textbf{0.01012} \\
\hline
Hybrid & Price-lookup & 252 & 42 & 28 & 0.1196 & 0.2392 \\
\hline
Hybrid & Attribute & 248 & 41 & 19 & \textbf{0.00622} & \textbf{0.02487} \\
\hline
\end{tabularx}
\end{table}
\FloatBarrier

After Holm correction: dense-attribute remains significant
($p_{\text{Holm}}=0.00438$); hybrid-attribute remains significant
($p_{\text{Holm}}=0.02487$); hybrid-full-set remains significant
($p_{\text{Holm}}=0.01012$); dense-full-set is no longer significant
($p_{\text{Holm}}=0.14455$); and neither price-lookup comparison is
significant. This is direct statistical evidence for the chapter's central
claim: the difference between table-aware and recursive chunking
concentrates in attribute questions; there is no statistical evidence that
table-aware chunking is better for price-lookup questions, and the
overall dense advantage does not survive correction for multiple
comparisons. This answers RQ4: with \texttt{gpt-oss-20b}, the table-aware
pipeline delivers a stable improvement on attribute questions in both dense
and hybrid retrieval, and in the hybrid configuration it is also better
overall after Holm correction.

\paragraph{1.2.8 Retrieval-to-answer conversion and context-use mechanism (RQ5).}
Table~\ref{tab:res-conversion} reports the share of retrieval hits converted
into correct answers.

\begin{table}[htbp]
\centering
\caption{Conversion rate: retrieval hit $\rightarrow$ correct answer}
\label{tab:res-conversion}
\small
\renewcommand{\arraystretch}{1.25}
\begin{tabularx}{\textwidth}{|X|>{\centering\arraybackslash}p{0.18\textwidth}|>{\centering\arraybackslash}p{0.18\textwidth}|>{\centering\arraybackslash}p{0.18\textwidth}|}
\hline
\rowcolor{tableheader}
\textbf{Config.} & \textbf{Overall} & \textbf{Price-lookup} &
\textbf{Attribute} \\
\hline
T1500 dense & 241/343 = 70.26\% & 159/189 = 84.13\% & \textbf{82/154 = 53.25\%} \\
\hline
T1500 hybrid & 287/410 = 70.00\% & 196/238 = 82.35\% & \textbf{91/172 = 52.91\%} \\
\hline
R1500 dense & 221/354 = 62.43\% & 162/196 = 82.65\% & 59/158 = 37.34\% \\
\hline
R1500 hybrid & 249/403 = 61.79\% & 182/228 = 79.82\% & 67/175 = 38.29\% \\
\hline
\end{tabularx}
\end{table}
\FloatBarrier

On price-lookup questions, T and R conversion rates sit close together
(79.8\%--84.1\%). On attribute questions the gap is substantially larger:
53.25\% versus 37.34\% (dense) and 52.91\% versus 38.29\% (hybrid).
Notably, T1500 hybrid's attribute Recall@5 is \emph{lower} than R1500's
(69.4\% versus 70.6\%) yet its generation accuracy is \emph{higher} (38.7\%
versus 29.8\%), which rules out the simple explanation that T answers better
purely because retrieval is better. BM25 raises Recall@5 sharply but barely
moves conversion (T1500: 70.26\%$\rightarrow$70.00\%; R1500:
62.43\%$\rightarrow$61.79\%) -- BM25 mainly adds more target values into the
top-5 without automatically making the model use context more effectively,
whereas table-aware chunking creates no attribute-retrieval advantage yet
delivers substantially higher attribute conversion. This answers RQ5: the
data is consistent with the hypothesis that table-aware chunking helps the
model exploit row--column relationships in context more effectively -- this
is indirect evidence from conversion, not an absolute causal proof, since
retrieval hits still rest on an answer-containing proxy and the two
branches' top-5 sets are not fully identical.

\paragraph{1.2.9 Representation effect by model, fixed retrieval.}
This experiment isolates the effect of \textbf{generation-time
representation} from the effect of retrieval: three arms (HTML, KV, VERB)
share the same T1500-dense top-5, chunk order, and cell data, changing only
the context string; a fourth arm (Recursive) uses R1500-dense top-5 and is
treated as a separate baseline pipeline, not a fourth representation sharing
T1500 chunks. Table~\ref{tab:res-by-model} summarizes each model's best
table-aware representation against the recursive baseline.

\begin{table}[htbp]
\centering
\caption{Best table-aware representation versus recursive baseline, by model (overall EM)}
\label{tab:res-by-model}
\small
\renewcommand{\arraystretch}{1.25}
\begin{tabularx}{\textwidth}{|X|X|>{\centering\arraybackslash}p{0.16\textwidth}|>{\centering\arraybackslash}p{0.14\textwidth}|}
\hline
\rowcolor{tableheader}
\textbf{Model} & \textbf{Best table-aware arm} & \textbf{Recursive} &
\textbf{Difference} \\
\hline
GPT-OSS 20B & HTML: 251/500 & 227/500 & +24 \\
\hline
Gemma 3 12B & HTML: 293/500 & 251/500 & +42 \\
\hline
Llama 3.1 8B & KV: 231/500 & 204/500 & +27 \\
\hline
Gemma 3 4B & HTML: 235/500 & 183/500 & +52 \\
\hline
\end{tabularx}
\end{table}
\FloatBarrier

Every model's best table-aware representation exceeds the recursive baseline
on overall EM. The pattern across representations, however, is not
monotonic with model size. For GPT-OSS~20B, HTML reaches the highest overall
EM (251/500), ahead of KV by 8 questions and VERB by 20; compared with KV,
HTML is 10 questions ahead on price-lookup but 2 behind on attributes. For
Gemma~3~12B, HTML leads on both overall and attribute EM, with a relatively
small gap between the three table-aware representations, especially on
attributes. For Llama~3.1~8B, KV is highest overall, with most of the gain
over HTML concentrated in price-lookup ($+19$) and only $+2$ on attributes.
For Gemma~3~4B, HTML exceeds both KV and VERB on both question groups (41
and 20 questions ahead of KV on price-lookup and attributes respectively;
34 and 36 ahead of VERB). These differences are reported as observed values;
only differences confirmed by paired McNemar with Holm correction on raw,
per-\texttt{query\_id} results are described as statistically significant in
this study. The result supports a clear conclusion: \textbf{no single
representation is optimal for every model.} The effectiveness of HTML,
key--value, and verbalized prose depends on model family and question type,
not simply on nominal parameter count -- so the generation-time
representation should be chosen against the deployed model's own benchmark
rather than inferred from model size.

\paragraph{1.2.10 Rejection ability on the 70 unanswerable questions (RQ6).}
Table~\ref{tab:res-rejection} reports correct-rejection results for the two
hybrid configurations.

\begin{table}[htbp]
\centering
\caption{Correct-rejection results, two hybrid configurations}
\label{tab:res-rejection}
\small
\renewcommand{\arraystretch}{1.25}
\begin{tabularx}{\textwidth}{|X|>{\centering\arraybackslash}p{0.16\textwidth}|>{\centering\arraybackslash}p{0.12\textwidth}|>{\centering\arraybackslash}p{0.12\textwidth}|>{\centering\arraybackslash}p{0.12\textwidth}|}
\hline
\rowcolor{tableheader}
\textbf{Config.} & \textbf{Overall} & \textbf{N1} & \textbf{N3} & \textbf{N4} \\
\hline
T1500 hybrid & 44/70 = 62.9\% & 13/30 = 43.3\% & 19/20 = 95.0\% & 12/20 = 60.0\% \\
\hline
R1500 hybrid & 47/70 = 67.1\% & 11/30 = 36.7\% & 16/20 = 80.0\% & 20/20 = 100\% \\
\hline
\end{tabularx}
\end{table}
\FloatBarrier

Under the current rubric, T1500 fails to reject correctly on 26/70 questions
(37.1\%), and R1500 on 23/70 (32.9\%). Paired comparison: T correct/R
incorrect on 9 questions, T incorrect/R correct on 12 questions, McNemar
$p=0.6636$ -- no evidence of an overall difference between the two
pipelines. Group-level results show different failure modes: both remain
weak on N1 (a row exists but the requested attribute cell is empty); T is
higher on N3 (near-miss variants that do not exist); R is higher on N4
(publication periods absent from the corpus). These N1/N3/N4 differences are
used only to describe failure patterns, not treated as independently
confirmed conclusions, since each subgroup is small and was not corrected as
its own testing family. This answers the safety portion of RQ6: table-aware
chunking does not improve overall rejection ability; both pipelines need an
evidence-verification mechanism and an explicit existence check on the
underlying data before emitting a specific value -- which is exactly why the
pricing tool in Chapter~V performs its own explicit verification chain
rather than relying on chunking or retrieval to provide this safety
property.

\paragraph{1.2.11 Best configurations by evaluation scope.}
Because the experiments were designed to isolate each component, three
different ``best'' configurations are reported by scope rather than forced
into a single, never-directly-run global optimum: the best retrieval
representation is \textbf{HTML} (Recall@5 = 343/500 = 68.6\%, versus 306/500
for KV and 218/500 for VERB), used as the default indexing representation;
the best generation-time configuration with retrieval held fixed is
\textbf{T1500 dense top-5 + HTML context + Gemma~3~12B}, reaching EM =
293/500 = 58.6\% (159/252 = 63.1\% price-lookup, 134/248 = 54.0\% attribute)
-- one question higher than the 292/500 hybrid GPT-OSS configuration below,
but not a direct end-to-end comparison since this benchmark holds dense
retrieval fixed while the GPT-OSS configuration uses hybrid retrieval; and
the best \textbf{fully confirmed end-to-end hybrid} configuration is T1500 +
HTML tables with header repetition + \texttt{add\_table\_context} disabled +
dense embedding + BM25~Okapi + RRF ($k=60$) + top-5 + \texttt{openai/gpt-oss-20b},
reaching Recall@5 = 410/500 = 82.0\%, EM = 292/500 = 58.4\% (196/252 = 77.8\%
price-lookup, 96/248 = 38.7\% attribute), and correct rejection = 44/70 =
62.9\%. These three results do not contradict one another: HTML is the
clearest retrieval choice, Gemma~3~12B~+~HTML is the strongest generation
choice under fixed evidence, and T1500~HTML~hybrid~+~GPT-OSS~20B is the
fully confirmed end-to-end pipeline. No further cross-optimization across
every model~$\times$~retriever~$\times$~representation combination was run,
since the primary goal was to identify mechanism and per-layer contribution
rather than locate a maximum over a large configuration space.

\subsection{2. Discussion}

\subsubsection{2.1 ASR Discussion}
\guidance{To be completed by the Speech AI Lead: discuss how ASR results
address the ASR research questions, how RouteLoRA's accent-routing behavior
compares with single-accent adapters and the PhoWhisper baseline, and any
error patterns observed by accent, noise condition, or domain terminology.}

\subsubsection{2.2 Agentic RAG Discussion}

\paragraph{Why table-aware chunking wins at generation but not at retrieval
in general.}
The four result layers form a consistent chain: T1500 preserves headers and
column relationships better than R1500; repeating headers makes the T1500
index carry more near-duplicate vectors; R1500 has higher R@1, R@3, and MRR
in dense retrieval, and remains higher on these metrics in hybrid retrieval;
at the top-5 operating point, T1500 hybrid reaches 410/500 versus R1500's
403/500, though R1500 still has higher attribute Recall@5 specifically
(175 versus 172); and despite that, T1500 shows markedly higher attribute
generation and conversion. The table-aware pipeline's main contribution is
therefore not a better retrieval ranking at every depth -- its benefit
appears once evidence is already inside the context window and the model
must correctly interpret row--column relationships. Price-lookup questions
depend less on column structure, so the two branches do not differ with
statistical significance; attribute questions require mapping a value to
the correct column, so T1500's advantage survives Holm correction.

\paragraph{The complementary roles of BM25 and table-aware chunking.}
BM25 substantially raises retrieval for both pipelines, especially on
price-lookup questions rich in identifier tokens (codes, sizes, proper
names, numeric strings) well suited to lexical matching. Table-aware
chunking solves a different problem: retaining column meaning. The two
components are therefore complementary -- BM25 improves the chance of
placing a relevant item in the top-5, while table-aware chunking improves
the ability to use row--column relationships within the retrieved chunks.
Because neither the retrieval interaction nor the generation interaction
reaches statistical significance, it would be inaccurate to say BM25 is
effective only in combination with table-aware chunking; the more accurate
conclusion is that BM25 helps generally, while table-aware chunking creates
its own, separate benefit concentrated in the attribute group.

\paragraph{Retrieval representation and generation representation must be
distinguished.}
The representation experiment shows HTML is not merely a convenient storage
format. When the same 2,252 logical T1500 chunks are embedded under three
representations, Recall@5 is 343 for HTML, 306 for KV, and 218 for
verbalized -- so \textbf{HTML is the best observed retrieval
representation} at the benchmark's primary endpoint. The fixed-evidence
generation benchmark shows HTML is also a strong generation-stage
representation: GPT-OSS~20B, Gemma~3~12B, and Gemma~3~4B all reach their
highest overall EM with HTML, while Llama~3.1~8B reaches its highest with
KV. The Llama case shows generation representation still depends on the
model, even though HTML wins for three of the four models evaluated.
Retrieval and generation therefore continue to need distinct objectives:
retrieval needs a string that embeds with strong discriminability between
chunks; generation needs a string that the specific deployed model can
exploit accurately in context. The architecture accordingly keeps the
\textbf{canonical table grid plus HTML} as the source of truth and the
default indexing representation, while still allowing deterministic
rendering to KV or verbalized prose once a deployed generation model's own
benchmark shows a benefit. There is no evidence for a simple rule such as
``the smaller the model, the more it needs KV or verbalized text.''

\paragraph{Implications for deployment.}
The results show that a RAG pipeline for price tables should not be
optimized against a single metric. If the primary goal is finding a code or
numeric value, hybrid retrieval has a large effect. If the goal is
answering attribute questions by column, preserving structure has a clear
value. If a document does not contain the answer, both pipelines still show
a moderate rejection rate, particularly for empty-cell cases -- which is
exactly why the deployed system (Chapter~V, Section~3.2.3) adds an explicit
verification layer before answering, such as checking whether a value
appears in the correct row, column, and data period, rather than relying on
chunking or retrieval alone to signal ``not found.'' A structured-data layer
or verification API is a reasonable direction for further evaluation, but
this study did not run a direct RAG-versus-SQL benchmark and therefore does
not conclude that one architecture is generally superior.

\paragraph{Connecting experimental results to architecture decisions.}
The six research questions posed in Chapter~IV operate at the mechanism
level -- structure, embedding, retrieval, generation, how the model uses
context, and reliability under rejection. Answering them does not by itself
produce an architecture; it identifies \emph{which component is responsible
for which benefit} and \emph{what that benefit costs}. Three conclusions
connect directly to the design decisions in Chapter~V. First, the benefit of
table-aware chunking is not winning retrieval in general, but appearing
specifically in the question group that requires mapping a value to a
column (RQ1, RQ4) and in the model's ability to exploit context once
evidence has been retrieved (RQ5) -- this boundary is the direct reason
Chapter~V splits the system into two complementary data paths (table-aware
RAG and a typed tool with structured queries) rather than using one RAG
mechanism for every question type. Second, the retrieval-layer
representation experiment (RQ3) elevates HTML from a ``convenient storage
format'' to a quantitatively supported choice -- on the same logical chunks,
HTML wins over KV and verbalized prose precisely at the Recall@5 operating
point -- direct enough evidence to feed straight into the architecture as
the default indexing representation. Third, the multi-model generation
benchmark shows no representation wins absolutely for every model, which
does not yield one fixed choice to hard-code, but instead yields a
\textbf{design constraint}: the representation used for a model to read
context must be kept as a per-model configuration parameter, decoupled from
the representation used for indexing.

\subsubsection{2.3 Integration Discussion}

The Agentic RAG results above assume a correctly transcribed text question.
In the full voice pipeline, a customer's spoken question first passes
through the ASR component (Chapter~V, Section~2) before reaching the router
described in Chapter~V, Section~3.2.4. \guidance{To be completed jointly
once ASR evaluation is available: discuss how ASR transcription errors are
expected to propagate into (a)~the router's keyword-based and
classifier-based decisions, (b)~material-name matching (Chapter~V, Section
3.2.3), which already required a dedicated relaxation algorithm to handle
even correctly-typed shortened names, and (c)~retrieval quality; and discuss
whether a low ASR-confidence transcript should route directly to
\texttt{CLARIFY} rather than being answered as if fully reliable.} Because
material names in this domain are dense with codes, sizes, and
manufacturer-specific terms -- exactly the token classes shown in Section~1.2.5
to be weak points for semantic retrieval and highly sensitive to exact
matching -- ASR misrecognition of a product code or size is expected to be a
disproportionately costly error mode compared with misrecognition of a
common word, and this should be validated once end-to-end voice evaluation
is available.

\subsection{3. Recommendations}

For the \textbf{Agentic RAG component}, four directions follow directly from
the results above: (1)~add an evidence-level gold label (document, table,
row, and column identifiers) for at least a subset of questions, to measure
evidence recall rather than the current answer-containing proxy, since
Section~1.2.3 already shows the proxy can both overstate and understate a
true retrieval hit; (2)~extend the rejection benchmark and the underlying
existence-verification chain, since Section~1.2.10 shows chunking strategy
alone does not reliably improve refusal on unanswerable questions;
(3)~evaluate a lightweight document-quality classifier as a fallback
trigger for the geometry-based extraction pipeline, activating a stronger
visual model only when simple quality signals (cell-fill rate, cross-page
column-count consistency) fall below a threshold, extending the same
conditional-activation philosophy already used for the OCR fallback path
(Chapter~V, Section~3.2.1) to complex-layout PDFs that still carry a text
layer; and (4)~re-evaluate the typed-tool-versus-text-to-SQL decision
(Chapter~V, Section~3.2.3) if the pricing schema grows to include linked
tables and multi-period price history, since the current justification is
explicitly scoped to the present schema size, not a permanent rule.

For the \textbf{ASR component}, \guidance{the Speech AI Lead should add
recommendations here, for example: expanding regional-accent data coverage
for underrepresented regions identified during error analysis; evaluating
RouteLoRA's routing-confidence calibration as a trigger for the
\texttt{CLARIFY} route described in Chapter~V; and re-running the price
material-name matching evaluation (Chapter~V, Section 3.2.3) using ASR
transcripts rather than typed text, to measure whether accent-specific
recognition errors interact with the token-relaxation safety gates already
tuned for typed queries.}

\subsection{4. Personal Reflection (Optional)}
\guidance{Reflect personally on the project process, lessons learned, and how
the work could be improved in the future.}

% -----------------------------------------------------------------------------
% VII. Conclusion
% -----------------------------------------------------------------------------
\chapterstart{VII. Conclusion}

\subsection{1. Summary of Findings}

\textbf{ASR track.} \guidance{To be completed by the Speech AI Lead once ASR
evaluation is finished: summarize overall and per-accent WER/CER for the
PhoWhisper baseline and RouteLoRA, and state whether the ASR research
questions posed in Chapter~IV were answered.}

\textbf{Agentic RAG track.} This track evaluated the table-aware pipeline
across the full chain from structure, embedding, and retrieval through to
generation, representation by model, and rejection ability. The results
show that the benefit of table-aware chunking is not a uniform improvement
at every layer. At the structure layer, T1500 preserves headers and
row--column relationships better than R1500, while row integrity is nearly
equivalent once label bias is corrected. At the embedding layer, T1500
carries more near-duplicate vector pairs than R1500, reflecting the cost of
header repetition. The multi-cutoff retrieval profile sharpens the
conclusion further: in dense retrieval, R1500 exceeds T1500 at R@1 (45.6\%
versus 41.4\%), R@3 (65.4\% versus 61.4\%), R@5 (70.8\% versus 68.6\%), and
MRR (0.5659 versus 0.5345), with T1500 only slightly ahead at R@10; in
hybrid retrieval, R1500 remains higher at R@1, R@3, R@10, and MRR, while
T1500 reaches a higher R@5 at exactly the system's top-5 operating point
(82.0\% versus 80.6\%). Table-aware chunking is therefore not described as
generally better at retrieval. The representation experiment on the same
2,252 logical chunks gives a clear result: HTML reaches Recall@5 =
343/500 (68.6\%), KV 306/500 (61.2\%), and verbalized 218/500 (43.6\%),
confirming HTML as the best observed dense-retrieval representation at the
primary endpoint -- linearizing into KV or prose does not improve
retrieval.

Table-aware chunking's strongest benefit appears in attribute generation:
T1500 dense reaches 251/500 versus R1500's 227/500; T1500 hybrid reaches
292/500 versus R1500's 256/500. After Holm correction, the attribute
difference remains significant in both dense ($p_{\text{Holm}}=0.00438$)
and hybrid ($p_{\text{Holm}}=0.02487$) retrieval, while neither price-lookup
comparison is significant. Conversion analysis reinforces this
interpretation: on attribute questions, T1500 converts a proxy retrieval hit
into a correct answer roughly 53\% of the time, versus roughly 37--38\% for
R1500, even where R1500's attribute Recall@5 is not lower. BM25 improves
retrieval and generation for both pipelines, mainly on price-lookup
questions, with no statistically significant interaction with chunking
strategy at either layer -- so BM25 is treated as a general lexical-retrieval
improvement, and table-aware chunking as the component that preserves
structural signal for the generation model. The multi-model benchmark
further shows no absolute winner in generation-time representation across
models, though HTML leads in three of four cases, rejecting a simple
``smaller models need KV or verbalized text'' rule. On the 70 unanswerable
questions, T1500 hybrid reaches 62.9\% correct rejection and R1500 hybrid
67.1\%, with no significant overall difference (McNemar $p=0.6636$) --
optimizing retrieval and EM alone is not sufficient to guarantee reliability
when the requested data does not exist.

Final results are reported by scope rather than forced into a single global
optimum: \textbf{HTML is the best observed retrieval representation at
Recall@5; Gemma~3~12B + HTML is the best generation configuration with
evidence held fixed; and T1500 HTML + dense/BM25 RRF + GPT-OSS~20B is the
fully confirmed end-to-end hybrid pipeline, reaching Recall@5 = 82.0\% and
EM = 58.4\%.} Overall: with long, multi-page tabular documents, table-aware
chunking does not necessarily rank evidence higher than recursive chunking
at every retrieval depth; its principal value is preserving column labels
so the model uses evidence more accurately on attribute questions. HTML is
best suited to dense indexing on this benchmark; BM25 adds complementary
lexical search capability; and generation-time representation must be
chosen per deployed model.

\subsection{2. Contributions and Reflections on the Project}

The project's central architectural contribution is not choosing between
RAG and structured querying, but separating the problem according to the
true nature of each information type: \textbf{table-aware RAG} handles
ambiguous questions, row--column relationships, standards, manufacturers,
specifications, and conditions of applicability; a \textbf{typed tool
combined with structured queries} handles unit prices, quantities, and
calculations that require deterministic precision; \textbf{query-time
orchestration} combines both paths for mixed questions and ensures the
``does this need a number'' decision is made before any RAG context can
interfere; and a \textbf{shared canonical extraction layer} guarantees the
RAG and SQL paths never see two different versions of the same table.

The Chapter~VI retrieval evidence makes this architectural decision more
precise than intuition alone would allow. T1500 does not beat R1500 at
every retrieval depth: R1500 has higher top-rank precision and MRR, while
T1500 hybrid only pulls ahead at Recall@5 -- the actual operating point the
deployed system uses. The architecture therefore does not rest on a blanket
claim such as ``table-aware retrieval is better,'' but on two separate,
independently measured benefits: BM25 improves lexical retrieval generally,
while table-aware chunking helps the model understand and use structure
correctly once evidence has been retrieved. The indexing representation was
also confirmed directly by data: on the same logical chunk set, HTML
clearly outperforms the two alternative formats at Recall@5, making it the
default representation for embedding and indexing. At the generation layer,
no format wins absolutely for every model tried, so the representation used
for a model to read context remains a per-model configuration parameter,
distinct from the indexing representation, which does not vary with the
generation model. The model actually running in production is likewise an
independent operational decision from the model used to run the generation
benchmark -- a distinction necessary to avoid confusing ``the configuration
that was measured'' with ``the configuration currently serving real users''
(Chapter~V, Section~3.1).

The project chose not to continue cross-optimizing across the full
model~$\times$~retriever~$\times$~representation combination space.
Instead, it stopped at conclusions measured independently per layer --
structure, embedding, retrieval, generation -- keeping each one separate
from operational decisions such as the choice of production model. This
trades some of the theoretical maximum achievable over a large configuration
space for an architecture that is easier to explain, easier to reproduce,
and better suited to the auditability an enterprise system requires, where
the ability to trace a number back to its exact source matters as much as
average accuracy on a benchmark.

\textbf{ASR track contributions.} \guidance{The Speech AI Lead should
summarize the specific technical contribution of the PhoWhisper +
RouteLoRA multi-accent recognition work once evaluation is complete.}

\guidance{Optional: add a brief personal reflection on what worked well,
what was most challenging, and what the team learned during development.}

\subsection{3. Limitations and Future Work}

\subsubsection{3.1 ASR Limitations and Future Work}
\guidance{To be completed by the Speech AI Lead. Suggested structure:
limitations of regional-accent data coverage, generalization beyond the
evaluated accents/noise conditions, and reproducibility constraints tied to
model/checkpoint versioning; proposed future work such as expanding
under-represented accent data, evaluating routing-confidence calibration,
and joint ASR+RAG end-to-end evaluation.}

\subsubsection{3.2 Agentic RAG Limitations and Future Work}

\textbf{Limitations.} (1)~The experimental corpus is limited to the
Vietnamese construction-material price-bulletin domain; benchmark figures do
not automatically generalize to other data domains. (2)~One very long table
document dominates a large share of chunks in the evaluation set, which may
affect sample representativeness. (3)~Retrieval-hit labels rely on whether
the expected answer appears in a chunk, not yet an independent gold label at
the row/column level. (4)~Attribute-question labels originate from the same
parser used to extract the data; manual verification on a small subsample
is a sanity check, not a full independent gold set. (5)~``Best''
configurations were selected on the same benchmark used to evaluate them,
without an independent held-out set. (6)~Generation results depend on the
specific model and serving version at measurement time; reproduction
requires locking the model, prompt, parameters, and a run-manifest file
(Chapter~V, Section~3.4). (7)~Scanned PDFs and tables without a text layer
were evaluated only to a limited extent through the OCR fallback path, with
no dedicated quantitative benchmark for this document type. (8)~Rejection
ability on unsupported questions remains moderate, particularly when the
source field is empty or the request falls outside the ingested data.
(9)~The full model~$\times$~retriever~$\times$~representation combination
space was not evaluated; conclusions are decomposed per layer to prioritize
explainability over locating a maximum over a large configuration space.
(10)~The pricing path does not yet support aggregate queries (minimum,
maximum, median by unit) -- a question like ``roughly what does cement
cost'' is not yet answered with a single aggregate figure (Chapter~V,
Section~3.2.3). (11)~For OCR'd tables with multiple parallel regional price
columns, the column-detection mechanism currently maps only one shared price
column, risking missed rows whose value sits only in a secondary column.
(12)~A portion of price data ingested before the unit-symbol-repetition
resolver existed still carries unnormalized values and needs a re-ingestion
or backfill pass to fully clean.

\textbf{Representation by layer versus by model -- two complementary, not
contradictory, conclusions.} The two Chapter~VI representation benchmarks
support, rather than contradict, each other, and this is the basis for
separating representation into three layers in the architecture. At the
\emph{retrieval} layer, holding the same T1500 logical chunks and embedding
them under three formats, HTML leads clearly at Recall@5
(Section~VI.1.2.4) -- hence HTML is retained as the default indexing
format. At the \emph{generation} layer, holding retrieved context fixed and
varying only the format the model reads, no format wins absolutely for
every model tried -- some models read HTML best, others read key--value
better -- which does not support a simple rule such as ``the smaller the
model, the simpler the format it needs.'' The architecture therefore
separates three independent representation layers: \emph{canonical data}
(the auditable row--column grid and HTML, Chapter~V Section~3.2.1),
\emph{retrieval representation} (HTML by default, locked in by Recall@5
evidence, independent of the generation model), and \emph{generation-time
representation} (renderable into another format from the same canonical
data depending on the deployed model, with HTML a reasonable default since
it leads for most models tested, but kept as a configuration parameter
rather than hard-coded). This separation means the system never needs to
re-parse a PDF from scratch when the generation model changes -- format
conversion happens only from the single canonical source and must preserve
headers, column order, empty cells, and every cell value.

\textbf{Future work.} (1)~\emph{Gold evidence and cross-document
benchmarking}: add a finer-grained gold label per question -- document,
table, row, and column identifiers containing the correct answer -- enabling
evidence recall measured against the true row rather than only whether the
expected answer string appears in a retrieved chunk; this distinction
matters when the same value appears at multiple positions in the corpus,
since only evidence-level measurement can confirm retrieval found the
\emph{correct} row rather than an incidentally matching value elsewhere.
(2)~\emph{Expanding to Text-to-SQL}: the current typed-tool decision suits
the present schema scale, not a permanent rule; text-to-SQL becomes worth
considering once several signals appear together -- multiple linked tables,
multi-period price history requiring time comparison, open-ended aggregate
queries that cannot be enumerated in advance (e.g., comparing price
movement across manufacturers over several quarters), and a question space
that exceeds what a fixed parameter schema can cover; adopting it would
require a separate read-only database, an explicit allow-list of accessible
tables/columns, an execution time limit, and a query-syntax validation step
before execution -- trading back part of the security surface gained from
typed tool-calling (Chapter~V, Section~3.4) for flexibility.
(3)~\emph{A quality-classification tier as a parser fallback}: extend the
CPU-first PDF extraction architecture (Chapter~V, Sections~3.2.1 and~3.3)
with a quality-classification tier -- run the default geometry-based path
first, check simple quality criteria on the result (e.g., readable-cell
rate, cross-page column-count consistency), and escalate to a stronger
visual document-understanding model only when those criteria are not met,
mirroring how the OCR fallback is already conditionally triggered, but
extended to text-layer PDFs whose table layout is too complex for the
default geometry path. (4)~\emph{Aggregate query support} on the pricing
path (minimum/maximum/median by unit), directly addressing Limitation~(10)
above. (5)~\emph{Joint ASR+RAG evaluation}, extending
Section~VI.2.3's integration discussion into a measured benchmark once ASR
transcripts are available, to quantify how transcription errors on product
codes and material names propagate into retrieval and pricing-tool
accuracy.

% -----------------------------------------------------------------------------
% VIII. References
% -----------------------------------------------------------------------------
\chapterstart{VIII. References}

\begingroup
\renewcommand{\section}[2]{}
\begin{thebibliography}{99}

\bibitem{gans2003}
N. Gans, G. Koole, and A. Mandelbaum,
``Telephone call centers: Tutorial, review, and research prospects,''
\textit{Manufacturing \& Service Operations Management}, vol. 5, no. 2,
pp. 79--141, 2003, doi: 10.1287/msom.5.2.79.16071.

\bibitem{ahmed1995}
S. M. Ahmed and R. Kangari,
``Analysis of client-satisfaction factors in construction industry,''
\textit{Journal of Management in Engineering}, vol. 11, no. 2,
pp. 36--44, 1995, doi: 10.1061/(ASCE)0742-597X(1995)11:2(36).

\bibitem{saka2023}
A. B. Saka, L. O. Oyedele, L. A. Akanbi, S. A. Ganiyu, D. W. M. Chan,
and S. A. Bello,
``Conversational artificial intelligence in the AEC industry: A review of
present status, challenges and opportunities,''
\textit{Advanced Engineering Informatics}, vol. 55, Art. no. 101869, 2023,
doi: 10.1016/j.aei.2022.101869.

\bibitem{nabavi2023}
A. Nabavi, I. Ramaji, N. Sadeghi, and A. Anderson,
``Leveraging natural language processing for automated information inquiry
from building information models,''
\textit{Journal of Information Technology in Construction}, vol. 28,
pp. 266--285, 2023, doi: 10.36680/j.itcon.2023.013.

\bibitem{huang2018}
M.-H. Huang and R. T. Rust,
``Artificial intelligence in service,''
\textit{Journal of Service Research}, vol. 21, no. 2, pp. 155--172, 2018,
doi: 10.1177/1094670517752459.

\bibitem{adam2021}
M. Adam, M. Wessel, and A. Benlian,
``AI-based chatbots in customer service and their effects on user
compliance,'' \textit{Electronic Markets}, vol. 31, no. 2, pp. 427--445,
2021, doi: 10.1007/s12525-020-00414-7.

\bibitem{mashaabi2022}
M. Mashaabi, A. Alotaibi, H. Qudaih, R. Alnashwan, and H. Al-Khalifa,
``Natural language processing in customer service: A systematic review,''
\textit{arXiv preprint arXiv:2212.09523}, 2022,
doi: 10.48550/arXiv.2212.09523.

\bibitem{suhaili2021}
S. M. Suhaili, N. Salim, and M. N. Jambli,
``Service chatbots: A systematic review,''
\textit{Expert Systems with Applications}, vol. 184, Art. no. 115461, 2021,
doi: 10.1016/j.eswa.2021.115461.

\bibitem{lewis2020}
P. Lewis et al.,
``Retrieval-Augmented Generation for knowledge-intensive NLP tasks,''
in \textit{Advances in Neural Information Processing Systems}, vol. 33,
pp. 9459--9474, 2020.

\bibitem{xu2024}
Z. Xu, M. J. Cruz, M. Guevara, T. Wang, M. Deshpande, X. Wang, and Z. Li,
``Retrieval-Augmented Generation with knowledge graphs for customer service
question answering,'' in \textit{Proceedings of the 47th International ACM
SIGIR Conference on Research and Development in Information Retrieval},
pp. 2905--2909, 2024, doi: 10.1145/3626772.3661370.

\bibitem{weizenbaum1966}
J. Weizenbaum,
``ELIZA---A computer program for the study of natural language communication
between man and machine,'' \textit{Communications of the ACM}, vol. 9, no. 1,
pp. 36--45, 1966, doi: 10.1145/365153.365168.

\bibitem{shin2021}
S. Shin and R. R. A. Issa,
``BIMASR: Framework for voice-based BIM information retrieval,''
\textit{Journal of Construction Engineering and Management}, vol. 147,
no. 10, Art. no. 04021124, 2021,
doi: 10.1061/(ASCE)CO.1943-7862.0002138.

\bibitem{lin2023}
W. Y. Lin,
``Prototyping a chatbot for site managers using Building Information Modeling
(BIM) and Natural Language Understanding (NLU) techniques,''
\textit{Sensors}, vol. 23, no. 6, Art. no. 2942, 2023,
doi: 10.3390/s23062942.

\bibitem{zheng2023}
J. Zheng and M. Fischer,
``Dynamic prompt-based virtual assistant framework for BIM information
search,'' \textit{Automation in Construction}, vol. 155, Art. no. 105067,
2023, doi: 10.1016/j.autcon.2023.105067.

\bibitem{le2024}
T.-T. Le, L. T. Nguyen, and D. Q. Nguyen,
``PhoWhisper: Automatic speech recognition for Vietnamese,''
in \textit{Proceedings of the ICLR 2024 Tiny Papers Track}, 2024.

\bibitem{hu2022}
E. J. Hu, Y. Shen, P. Wallis, Z. Allen-Zhu, Y. Li, S. Wang, L. Wang,
and W. Chen,
``LoRA: Low-rank adaptation of large language models,''
in \textit{Proceedings of the International Conference on Learning
Representations}, 2022.

\bibitem{bagat2025}
R. Bagat, I. Illina, and E. Vincent,
``Mixture of LoRA experts for low-resourced multi-accent automatic speech
recognition,'' in \textit{Proceedings of Interspeech 2025}, pp. 1143--1147,
2025, doi: 10.21437/Interspeech.2025-775.

\bibitem{yao2023}
S. Yao, J. Zhao, D. Yu, N. Du, I. Shafran, K. Narasimhan, and Y. Cao,
``ReAct: Synergizing reasoning and acting in language models,''
in \textit{Proceedings of the International Conference on Learning
Representations}, 2023.

\end{thebibliography}
\endgroup

% -----------------------------------------------------------------------------
% IX. Appendices
% -----------------------------------------------------------------------------
\chapterstart{IX. Appendices}
\guidance{Provide the final appendices following the template as part II in
Report 9.}
\guidance{Include supplementary information that does not fit into the main
sections but remains important for a complete understanding of the project,
such as raw data, source code, full surveys, detailed diagrams, or other
supporting documents.}
\guidance{After completing all sections, ensure that the document is logically
and coherently organized and has been checked carefully for language,
spelling, and grammar errors.}

\end{document}
